\documentclass[11pt]{amsart}

\usepackage[T1]{fontenc}
\usepackage[utf8]{inputenc}
\usepackage{lmodern}
\usepackage{amsmath}
\usepackage{amssymb}
\usepackage{amsfonts}
\usepackage{amsthm}
\usepackage{mathtools}
\usepackage{geometry}
\usepackage{enumitem}
\usepackage[hidelinks]{hyperref}

\geometry{margin=1in}
\numberwithin{equation}{section}

\theoremstyle{plain}
\newtheorem{theorem}{Theorem}[section]
\newtheorem{proposition}[theorem]{Proposition}
\newtheorem{lemma}[theorem]{Lemma}
\newtheorem{corollary}[theorem]{Corollary}

\theoremstyle{definition}
\newtheorem{definition}[theorem]{Definition}
\newtheorem{hypothesis}[theorem]{Hypothesis}
\newtheorem{heuristic}[theorem]{Heuristic}
\newtheorem{question}[theorem]{Question}

\theoremstyle{remark}
\newtheorem{remark}[theorem]{Remark}

\DeclareMathOperator{\sgn}{sgn}
\DeclareMathOperator{\li}{li}
\newcommand{\Ms}{M_{\mathrm{s}}}
\newcommand{\Msk}{M_{\mathrm{sk}}}
\newcommand{\Ap}{\mathcal{A}_{+}}
\newcommand{\Am}{\mathcal{A}_{-}}
\newcommand{\Anet}{\mathcal{A}}
\newcommand{\eps}{\varepsilon}
\newcommand{\Pp}{P^{+}}
\newcommand{\Cc}{\mathcal{C}}
\newcommand{\Sq}{\mathcal{S}}
\newcommand{\En}{\mathcal{E}}
\newcommand{\mt}{\widetilde{m}}
\newcommand{\F}{\mathcal{F}}

\title[The Mertens excursion skeleton]{The area function of the Mertens excursion skeleton\\
and horizon-localized criteria for the Riemann hypothesis}

\author{V.~Casa\~na-Giner}
\date{}

\subjclass[2020]{Primary 11N37; Secondary 11M26, 60G70}
\keywords{Mertens function, M\"obius function, excursions, Riemann hypothesis,
Kaczorowski--Pintz oscillation, lacunary sequences, Dyck paths, Catalan tails, extreme values}

\begin{document}

\begin{abstract}
We decompose the Mertens function $M(x)=\sum_{n\le x}\mu(n)$ into its
excursions between zero runs and study a regularized triangular skeleton of
these excursions, its area curves, and exact quadratic and correlation
identities inside each completed excursion.

For any ordered family of disjoint integer intervals $(x_0,x_1]$, elementary
counting shows that the members starting at most $X$ with $x_1>x_0^A$
$(A>1)$ are $O_A(\log\log X)$, those with $x_1>cx_0^2$ are
$O_c(\log\log X)$, those with $x_1>x_0\psi(x_0)$ are $o(\log X)$ whenever
$\psi\to\infty$, and those with $x_1>\lambda x_0$ are at most
$\log X/\log\lambda+O(1)$. Combined with the Kaczorowski--Pintz oscillation
theorem, this gives a horizon-localization principle: if the excursions
starting at most $X$ that belong to an exceptional family or cross a horizon
$\mathfrak h$ number at most $(\kappa+o(1))\log X$ with $\kappa<\gamma_1/\pi$,
then the Riemann hypothesis is equivalent to Littlewood's bound
$|M(t)|\ll_\eps t^{1/2+\eps}$ restricted to $x_0<t\le\min\{x_1,\mathfrak h(x_0)\}$
on the remaining excursions. The horizon can be lowered from $x_0^2$ through
$cx_0^2$ and $x_0\psi(x_0)$ to $\lambda x_0$ for every
$\lambda>\lambda_*:=e^{\pi/\gamma_1}$, the exact threshold of this counting
mechanism. The same input shows that late excursions have length
$x_0^{1+o(1)}$. On the resulting window every $\theta>\theta_*:=\frac12e^{\pi/\gamma_1}$
is an admissible prefix-dependency cutoff: with $M(u)=0$ and
$u<t\le\lambda u<2u$,
\[
M(t)=-\sum_{d\le\lambda u/2}\mu(d)
\Bigl(\Bigl\lfloor\frac td\Bigr\rfloor-\Bigl\lfloor\frac ud\Bigr\rfloor\Bigr).
\]

After deleting zero increments, a completed excursion is a primitive Dyck
path on its fixed support mask. Under the uniform primitive Catalan
reference law on that mask, a value $|M(t)|\ge t^{1/2+\eps}$ in an excursion
with $x_1\le\lambda x_0$ has probability at most $x_0^{O(1)}e^{-x_0^{2\eps}/2}$,
and polynomial lower bounds for the complete-excursion or open-prefix tail
weights imply $|M(t)|\ll\sqrt{t\log t}$ on the corresponding ranges.

The one-step recursion concentrates, for every fixed $K\ge2$, all dependence
above $\lambda u/(K+1)$ in the $K-1$ channels $M(t/k)-M(u/k)$, $2\le k\le K$,
and a persistence lemma quantifies how a large value must be carried by one
of them or by the deeper remainder. A multiplicative-layer filtration shows
that for every fixed depth $k$ the first $k$ M\"obius layers determine the
$(k+1)$-st layer of the future window; the determined positions have density
zero, while their signed contribution is of order $u(\log\log u)^k/\log u$.
Thus substantial cancellation with the remaining layers is necessary to
reach the RH scale; no assertion on the number of layers required is made.

Under a uniform sub-Weibull tail hypothesis for a normalized excursion
statistic, the counting bounds give extreme-value scales by a union bound
and a Borel--Cantelli argument. We also give an exact bilinear
representation of $M$ on the square range of an excursion, an exact Mellin
decomposition of $1/\zeta$ into completed excursion blocks, and asymptotic
decompositions of $M$ into smooth and large-prime parts.
\end{abstract}

\maketitle

\tableofcontents

\section{Introduction}
\label{sec:introduction}

The object of this paper is the Mertens function $M$, decomposed into
excursions between zero runs (Section~\ref{sec:definitions}). Two kinds of
information about $M$ are separated throughout: what is forced by the
ordering and local geometry of the excursions alone, and what requires
arithmetic input.

\subsection*{Main results}

The central results are in Section~\ref{sec:lacunary}.

\smallskip\noindent\emph{Geometric counts.} For any ordered family of
disjoint integer intervals $(x_{0},x_{1}]$, the members with $x_1>x_0^A$
($A>1$) starting below $X$ are $O_A(\log\log X)$, those with
$x_1>cx_0^2$ are $O_c(\log\log X)$, those with $x_1>x_0\psi(x_0)$ are
$o(\log X)$ whenever $\psi\to\infty$, and those with $x_1>\lambda x_0$ are at
most $\log X/\log\lambda+O(1)$ (Lemmas~\ref{lem:power-lacunary},
\ref{lem:partition}, \ref{lem:superlinear}, \ref{lem:linear-lacunary}). No
arithmetic enters.

\smallskip\noindent\emph{Horizon localization.} By
Kaczorowski--Pintz~\cite{KP86}, if $\Theta>1/2$ then at least
$(\gamma_1/\pi-o(1))\log X$ distinct excursions starting below $X$ contain
values of size $t^{\Theta-\eta}$. Theorem~\ref{thm:horizon} turns every
geometric count below this threshold into an equivalent form of RH: for any
horizon $h$ whose crossing family has at most $(\kappa+o(1))\log X$ members
with $\kappa<\gamma_1/\pi$,
\[
\mathrm{RH}\iff |M(t)|\ll_\eps t^{1/2+\eps}\quad\text{for }x_0<t\le\min\{x_1,\mathfrak h(x_0)\}.
\]
In particular (Corollary~\ref{cor:descending})
\[
\mathrm{RH}\iff\text{Littlewood's bound holds on }x_0<t\le\min\{x_1,\lambda x_0\},
\qquad\lambda>\lambda_*:=e^{\pi/\gamma_1}=1.2488\ldots,
\]
and the square horizon $x_0^2$ is eroded from above through $cx_0^2$ and
$x_0\psi(x_0)$. If RH fails, a positive proportion (on the scale $\log X$) of the
Kaczorowski--Pintz excursions carry their large value in
$[x_0+\frac12x_0^{\Theta-\eta},\lambda x_0]$ (Corollary~\ref{cor:two-sided}). The threshold $\lambda_*$ is the multiplicative distance
between consecutive sign changes of $\cos(\gamma_1\log x)$
(Remark~\ref{rem:lambda-star}). It is the exact threshold of this counting mechanism, not a
proved limit for other localization methods. The local form of the
Kaczorowski--Pintz theorem also shows that for every fixed $A>1$ all late
excursions satisfy $x_1\le x_0^A$ (Corollary~\ref{cor:KP-local-returns}):
the length of a Mertens excursion is $x_0^{1+o(1)}$.  The fractional-square lemma makes the
intermediate erosion explicit:
\[
x_0^2\longrightarrow \frac{x_0^2}{2}\longrightarrow
\frac{x_0^2}{4}\longrightarrow\frac{x_0^2}{8}\longrightarrow\cdots
\longrightarrow x_0\psi(x_0)\longrightarrow\lambda x_0.
\]
For each fixed dyadic level the crossing family is $O(\log\log X)$; the
passage to a level depending on $x_0$ is supplied by the superlinear-horizon
lemma, not by a uniform use of the fixed-$c$ estimate.

Kaczorowski--Pintz~\cite{KP86} provide the oscillation-density input.  The
excursion-counting lemmas, the horizon-localization theorem, the
descending-horizon principle, the threshold
$\lambda_*=e^{\pi/\gamma_1}$, and the resulting prefix-dependency cutoff
$\theta_*=\tfrac12e^{\pi/\gamma_1}$ are developed here from that input.

\smallskip\noindent\emph{Linear window.} For $u<t\le2u$ one has
$\mu(n)=-\sum_{d\mid n,\,d\le u}\mu(d)$.  After horizon localization,
if $u<t\le\lambda u<2u$, the terms with $d>\lambda u/2$ have no multiple
in $(u,t]$, and therefore
\[
M(t)=-\sum_{d\le\lambda u/2}\mu(d)
\Bigl(\Bigl\lfloor\frac td\Bigr\rfloor-\Bigl\lfloor\frac ud\Bigr\rfloor\Bigr).
\]
Thus every $\theta>\theta_*:=\frac12e^{\pi/\gamma_1}$ is an admissible
prefix cutoff for the RH-equivalent excursion window; the convenient choice
$\lambda=5/4$ uses only $d\le5u/8$.  This is a descent of dependency, not
of size.  More sharply, for every fixed $K\ge2$ the dependence above
$\lambda u/(K+1)$ is concentrated in the $K-1$ increments
$M(t/k)-M(u/k)$, $2\le k\le K$; for $K=2$ and $\lambda=5/4$, the only channel
above $5u/12$ is the increment across $(u/2,t/2]$.  The one-step recursion
does not by itself give an iterable maximum principle, and a termwise
absolute-value contraction gives no nontrivial gain.  The exact
intra-excursion formulation, the scale cascade, a persistence extraction
lemma, and a ballistic counterexample to the relaxed divisor-sum version are
given below. A bounded number of multiplicative layers determines only a
density-zero set of increments in the localized window, but for every fixed
depth $k$ the signed contribution of the determined layers is of order
$u(\log\log u)^k/\log u$, far above the RH scale; reaching the RH scale
requires substantial cancellation with the remaining layers, and the present
arguments do not determine how many layers are needed.

\smallskip\noindent\emph{Extreme values.} If a normalized excursion statistic
has a uniform tail $\Pr(Z>z)\le Ce^{-\kappa_0z^\varpi}$, the counting bounds force the
maximum over a family with at most $\Lambda(X)$ members starting below $X$ to be
of order $(\log \Lambda(X))^{1/\varpi}$ in probability, with an almost sure $\limsup$ constant
$\kappa_0^{-1/\varpi}$ (Theorems~\ref{thm:ev-union} and~\ref{thm:ev-as}). Each level of
lacunarity removes one logarithm from the extreme scale:
$\sqrt{\log X}$, $\sqrt{\log\log X}$, $\sqrt{\log\log\log X}$.

\smallskip
The remaining sections contain exact identities for the skeleton and its area
curves (Sections~\ref{sec:definitions}--\ref{sec:formulas}), exact
correlation identities and sharp bounds inside an excursion
(Section~\ref{sec:internal_corr}), the relation with $\zeta$
(Section~\ref{sec:zeta}), area bounds and global criteria equivalent to RH
(Sections~\ref{sec:bounds}--\ref{sec:criteria}), the exact bilinear form of
$M$ on the square range of an excursion and the Mellin mass of an excursion
(Section~\ref{sec:lacunary}), and the decomposition of $M$ by large prime
factors, globally and inside an excursion
(Sections~\ref{sec:prime_horizon}--\ref{sec:local_large_prime}).
Section~\ref{sec:data} records numerical data on the excursion census.

\subsection*{Conventions}
$\rho$ runs over nontrivial zeros of $\zeta$, with ordinate $\gamma=\Im\rho$,
\[
\Theta=\sup\{\Re\rho\},\qquad \gamma_1=14.134725\ldots
\]
is the ordinate of the lowest zero, and RH means $\Theta=1/2$. Littlewood's
criterion \cite[\S14.25]{Titchmarsh1986} states that RH holds if and only if
$M(x)\ll_\eps x^{1/2+\eps}$ for every $\eps>0$. Statements labelled
\emph{Heuristic} are not used in any proof.

\subsection*{Notation}
\label{subsec:notation}
The following symbols keep one meaning throughout; symbols defined and used
only inside a single proof are not listed. The letters $c,C$ (with or without
subscripts other than those listed) denote positive constants whose value may
change from line to line.

\smallskip
\noindent\emph{Zeros and constants.}
$\rho$, $\gamma=\Im\rho$: nontrivial zeros of $\zeta$;
$\Theta=\sup\Re\rho$;
$\gamma_1=14.1347\ldots$;
$\lambda_*=e^{\pi/\gamma_1}=1.24889\ldots$, the threshold of the counting
mechanism of Theorem~\ref{thm:horizon};
$\theta_*=\lambda_*/2=0.62444\ldots$, the prefix cutoff;
$M_0=M+2$.

\smallskip
\noindent\emph{Excursions.}
$I_\nu=(x_{0,\nu},x_{1,\nu}]$ an excursion, also written $(u_\nu,v_\nu]$ or
$(u,v]$; $\sigma_\nu\in\{\pm1\}$ its orientation (always $\sigma$, never a
real part); $\ell_\nu=x_{1,\nu}-x_{0,\nu}$; $H_\nu$ its height; $\tau_\nu$ its
apex; $a_\nu=\tau_\nu-x_{0,\nu}$, $b_\nu=x_{1,\nu}-\tau_\nu$;
$u(x)$, $L_x$: last zero $\le x$ and length of the terminal block.

\smallskip
\noindent\emph{Skeleton and curves.}
$\Ms$, $\Msk$: regularized and unregularized skeletons;
$\Ap,\Am,\Anet$: area curves; $\widetilde A(x)=\int_0^xM$.

\smallskip
\noindent\emph{Quadratic data.}
$S_t$, $S_r$: oriented partial sums inside a block or excursion;
$C_h$, $C_h^{(\nu)}$: lag correlations; $q_\nu$: number of nonzero steps;
$Q_\nu$, $Q_{\rm term}$: quadratic masses; $\beta_\nu$, $\beta_{\rm term}$:
absolute masses $\sum|M|$; $P_h^{(\nu)},D_h^{(\nu)}$: support and
opposite-sign pair counts; $\En$: energy.

\smallskip
\noindent\emph{Families and counts.}
$\F$ a family of excursions; $E_\F(X)$ the number of its members starting at
most $X$; $\mathfrak h$ a horizon function and $\F_{\mathfrak h}$ the family
crossing it; $E_A$, $E_\square=E_2$, $E_\lambda$ the counts for
$x_1>x_0^A$, $x_1>x_0^2$, $x_1>\lambda x_0$; $c$ in $cx_0^2$ is a fixed
horizon fraction; $\kappa$ the density constant in \eqref{eq:horizon-hyp};
$N_\eta(Y)$ the Kaczorowski--Pintz excursion count.

\smallskip
\noindent\emph{Extreme values.}
$Z_k$ normalized excursion statistics; $Z_{\max,\F}(X)$ their family maximum;
$C,\kappa_0,\varpi$ the constants of $(\mathrm{SG}_\varpi)$;
$\Lambda(X)$ a deterministic bound for $E_\F(X)$.

\smallskip
\noindent\emph{Linear window and recursion.}
$\mu_u=\mu\mathbf 1_{[1,u]}$; $\mathcal B_u(t)$ the bilinear form;
$c_q(t,u)$ the kernel coefficients, with $e_q$ and $\mathrm{Err}(t,u)$
their remainder; $\Delta_k(u,t)=M(t/k)-M(u/k)$ the channels and
$\mathcal R_K(u,t)$ the deep remainder; $X_k,Y_K$ their oriented values;
$\alpha\in(0,1)$ the exponent in the normalized ratios
$|M(q)|/q^\alpha$; $w_r$, $K_q$ weights and combined kernel;
$\delta_d$ a relaxed prefix and $D_\delta(n)$ its divisor sums;
$\Omega,\omega$ the number of prime factors with and without multiplicity.

\smallskip
\noindent\emph{Catalan comparison.}
$\Sq(x)=\#\{n\le x:\mu(n)\ne0\}$ the squarefree-compressed clock;
$L^\#(I)$ the compressed length; $r_I(t)$ the support time;
$\mathrm{Cat}_n$ the Catalan numbers; $p_{\mathrm{Cat}}(I,t)$,
$p^{\max}_{\mathrm{Cat}}(I)$, $p_{\mathrm{open}}(t)$ reference tail weights;
$\mathcal I_{\mathrm{Cat}}$ selection energy; $\mathcal M_r$ positive meanders.

\smallskip
\noindent\emph{Large-prime decomposition.}
$U\ge2$ a level (not an excursion start); $M^{\rm sm}_U$, $M^{\rm lp}_U$ the
$U^2$-smooth and large-prime parts; $R=t/U^2$; $N_y(t,U)$ copy
multiplicities; $\Cc(R)$, $\Upsilon(R)$ the coherent and multiplicity-weighted
coefficients; $\mt(y)=\sum_{a\le y}\mu(a)/a$; $\Pp(n)$ the largest prime
factor; $\Pi_a(t)$ local prime counts.

\smallskip
\noindent\emph{Mellin.}
$s=\varsigma+it$ in Section~\ref{sec:criteria} and in
Proposition~\ref{prop:excursion-mellin}; $D_I(s)$ the excursion block of
$1/\zeta$; $J$ a leg in Proposition~\ref{prop:dipoles}.

\section{Excursions and the regularized skeleton}
\label{sec:definitions}

Let $\mu$ be the M\"obius function and
\begin{equation}
\label{eq:mertens_def}
M(x)=\sum_{n\le x}\mu(n),\qquad x\ge0,
\end{equation}
piecewise constant and right-continuous, with jumps $\mu(n)\in\{-1,0,1\}$
at the integers. The segmentation of $M$ by zero runs and the record of its
extrema are standard computational objects; Hurst~\cite{Hurst18} computed
$M(x)$ for all $x\le10^{16}$. The area and quadratic functionals of lattice
excursions are classical
\cite{Takacs91,NguyenThe04,Richard05,Janson07,DKW15,CP17}; here they are
applied to the deterministic M\"obius walk.

\begin{definition}[Excursions and zero runs]
\label{def:excursion}
An \emph{excursion} is an interval
\[
I_\nu=(x_{0,\nu},x_{1,\nu}]
\]
with integer endpoints, $M(x_{0,\nu})=M(x_{1,\nu})=0$,
$M(x_{0,\nu}+1)=\sigma_\nu\in\{\pm1\}$, and
$\sigma_\nu M(n)\ge1$ for every integer $x_{0,\nu}<n<x_{1,\nu}$. Thus
$x_{1,\nu}$ is the first return to zero after the departure from
$x_{0,\nu}$. Consecutive excursions may be separated by a \emph{zero run}
$[x_{1,\nu},x_{0,\nu+1}]$, a maximal set of consecutive integers on which
$M=0$; it may consist of a single point. The first excursion is $(0,2]$.
Since $M$ takes values of both signs of size $\gg\sqrt x$ infinitely often
(Remark~\ref{rem:omega}) and has unit steps, it has infinitely many zeros
and every excursion is finite. We also write $u_\nu=x_{0,\nu}$ and
$v_\nu=x_{1,\nu}$.

Put
\[
\ell_\nu=x_{1,\nu}-x_{0,\nu},\qquad
H_\nu=\max_{x_{0,\nu}<n<x_{1,\nu}}|M(n)|,
\]
\[
\tau_\nu=\min\{n\in I_\nu:|M(n)|=H_\nu\},\qquad
a_\nu=\tau_\nu-x_{0,\nu},\qquad b_\nu=x_{1,\nu}-\tau_\nu.
\]
Every excursion has $H_\nu\ge1$ and $\ell_\nu\ge2$.
\end{definition}

\begin{definition}[Completed excursions and open prefixes]
\label{def:completed-open}
An excursion in the sense of Definition~\ref{def:excursion} is always
\emph{completed}: both $M(x_{0,\nu})=0$ and $M(x_{1,\nu})=0$ are part of
the definition. A segment $(u,x]$ with $M(u)=0$ and $M(n)\ne0$ for
$u<n\le x$, with no return condition at $x$, is an \emph{open prefix}; the
terminal block of Definition~\ref{def:Ms-terminal} is one. Any estimate
whose proof uses the backward bound $|M(u+r)|\le x_{1,\nu}-u-r$ holds for
completed excursions only and is never applied to an open prefix.
\end{definition}

\begin{lemma}[Height and step count]
\label{lem:height_steps}
For every excursion,
\begin{equation}
\label{eq:height_steps}
H_\nu\le\min(a_\nu,b_\nu)\le\frac{\ell_\nu}{2}.
\end{equation}
\end{lemma}

\begin{proof}
The steps of $M$ are $0$ or $\pm1$, so at least $H_\nu$ steps are needed to
reach height $H_\nu$ from zero and at least $H_\nu$ more to return to zero.
\end{proof}

\begin{definition}[Crossing and touching zero runs]
\label{def:crossing_touching}
The zero run $[x_{1,\nu},x_{0,\nu+1}]$ between consecutive excursions is a
\emph{crossing run} if $\sigma_{\nu+1}=-\sigma_\nu$ and a \emph{touching run}
if $\sigma_{\nu+1}=\sigma_\nu$.
\end{definition}

\begin{definition}[Regularized skeleton]
\label{def:Ms}
The \emph{skeleton} $\Ms$ is the continuous, piecewise linear function
defined on every excursion by
\begin{equation}
\label{eq:Ms_def}
\Ms(t)=
\begin{cases}
\sigma_\nu\sqrt{H_\nu}\dfrac{t-x_{0,\nu}}{a_\nu},
 &x_{0,\nu}\le t\le\tau_\nu,\\[2ex]
\sigma_\nu\sqrt{H_\nu}\dfrac{x_{1,\nu}-t}{b_\nu},
 &\tau_\nu\le t\le x_{1,\nu},
\end{cases}
\end{equation}
and by $\Ms\equiv0$ on zero runs. The \emph{unregularized skeleton} $\Msk$
is defined in the same way with peak $\sigma_\nu H_\nu$.
\end{definition}

\begin{definition}[Terminal block of a truncated computation]
\label{def:Ms-terminal}
Let $N\ge2$ and $u=\max\{n\le N:M(n)=0\}$. If $u=N$ there is no terminal
block. Otherwise $M(n)\ne0$ for $u<n\le N$, and since $M$ has steps
$0,\pm1$ all these values have the sign of $M(N)$. Put
\[
\sigma_{\rm term}=\sgn M(N),\qquad
H_{\rm term}=\max_{u<n\le N}|M(n)|,\qquad
\tau_{\rm term}=\min\{n\in(u,N]:|M(n)|=H_{\rm term}\}.
\]
On $[u,N]$ the regularized skeleton is the segment from $(u,0)$ to
$(\tau_{\rm term},\sigma_{\rm term}\sqrt{H_{\rm term}})$ followed, if
$\tau_{\rm term}<N$, by the segment ending at
$(N,\sigma_{\rm term}\sqrt{|M(N)|})$. Its signed area is
\[
\Delta\Anet(u,N]=\frac{\sigma_{\rm term}}2
\Bigl[\sqrt{H_{\rm term}}(N-u)+
\sqrt{|M(N)|}(N-\tau_{\rm term})\Bigr].
\]
The terminal skeleton is provisional: after the next return to zero the
completed excursion acquires a different descending side. Independently of
it, for real $x$ with last zero $u(x)\le\lfloor x\rfloor$ define
\[
\beta_{\rm term}(x):=\sum_{u(x)<n\le\lfloor x\rfloor}|M(n)|,\qquad
Q_{\rm term}(x):=\sum_{u(x)<n\le\lfloor x\rfloor}M(n)^2,\qquad
L_x:=\lfloor x\rfloor-u(x).
\]
These depend only on $M$.
\end{definition}

\begin{definition}[The three curves]
\label{def:areas}
With $\Ms^{\pm}=\max(\pm\Ms,0)$,
\begin{equation}
\label{eq:areas_def}
\Ap(x)=\int_0^x\Ms^+(t)\,dt,\qquad
\Am(x)=-\int_0^x\Ms^-(t)\,dt,\qquad
\Anet(x)=\int_0^x\Ms(t)\,dt.
\end{equation}
$\Ap$ is nondecreasing, $\Am$ is nonincreasing, and both are constant on
zero runs and on excursions of the opposite sign.
\end{definition}

\begin{remark}[The family of $g$-skeletons]
\label{rem:g-skeleton}
For any $g:\mathbb N_0\to[0,\infty)$ with $g(0)=0$ and $g(h)>0$ for $h\ge1$,
one may replace the apex height $\sqrt H$ by $g(H)$. The signed area
increment becomes $\tfrac12\sigma_\nu g(H_\nu)\ell_\nu$. The correlation
identities of Section~\ref{sec:internal_corr} concern the M\"obius
increments and do not depend on $g$. The choice $g(H)=\sqrt H$ is a
normalization.
\end{remark}

\section{Exact formulas}
\label{sec:formulas}

\begin{lemma}[Increment per excursion]
\label{lem:jump}
For every complete excursion,
\begin{equation}
\label{eq:jump}
\Anet(x_{1,\nu})-\Anet(x_{0,\nu})
=\sigma_\nu\frac{\sqrt{H_\nu}\,\ell_\nu}{2},
\qquad
\sum_{n\in I_\nu}\Ms(n)
=\sigma_\nu\sqrt{H_\nu}\frac{\ell_\nu}{2}.
\end{equation}
\end{lemma}

\begin{proof}
The two pieces are triangles of height $\sqrt{H_\nu}$ and bases $a_\nu,b_\nu$.
For the discrete sum, the ascending values at $j=1,\ldots,a_\nu$ sum to
$\sqrt{H_\nu}\sum_{j\le a_\nu}j/a_\nu=\sqrt{H_\nu}(a_\nu+1)/2$; the
descending values at $\tau_\nu+r$, $r=1,\ldots,b_\nu$, sum to
$\sqrt{H_\nu}\sum_{r\le b_\nu}(b_\nu-r)/b_\nu=\sqrt{H_\nu}(b_\nu-1)/2$.
\end{proof}

\begin{proposition}[Explicit form of the three curves]
\label{prop:explicit}
Let $x\in I_\nu$. Then
\begin{equation}
\label{eq:explicit_A}
\Anet(x)=
\begin{cases}
\Anet(x_{0,\nu})+\sigma_\nu\sqrt{H_\nu}
\dfrac{(x-x_{0,\nu})^2}{2a_\nu},&x\le\tau_\nu,\\[2ex]
\Anet(x_{0,\nu})+\sigma_\nu\sqrt{H_\nu}
\left(\dfrac{a_\nu}{2}+(x-\tau_\nu)-
\dfrac{(x-\tau_\nu)^2}{2b_\nu}\right),&x\ge\tau_\nu.
\end{cases}
\end{equation}
At excursion endpoints,
\begin{equation}
\label{eq:explicit_net}
\Anet(x_{1,k})=\sum_{\nu\le k}\sigma_\nu\frac{\sqrt{H_\nu}\ell_\nu}{2},\qquad
\Ap(x_{1,k})=\sum_{\substack{\nu\le k\\\sigma_\nu=+1}}
\frac{\sqrt{H_\nu}\ell_\nu}{2},\qquad
\Am(x_{1,k})=-\sum_{\substack{\nu\le k\\\sigma_\nu=-1}}
\frac{\sqrt{H_\nu}\ell_\nu}{2}.
\end{equation}
The curves are piecewise quadratic and $C^1$.
\end{proposition}

\begin{proof}
Integrate \eqref{eq:Ms_def}; the endpoint formulas follow from
Lemma~\ref{lem:jump}.
\end{proof}

A point $y$ is a \emph{strict inflection point} of a function if there
exists $\delta>0$ such that the function is strictly convex on one side of
$y$ and strictly concave on the other, within distance $\delta$. The second
derivative need not exist at $y$.

\begin{lemma}[Critical geometry of $\Anet$]
\label{lem:critical}
Let $R=[p,q]$ be the zero run between an excursion of sign $\sigma_-$ ending
at $p$ and the next excursion of sign $\sigma_+$ starting at $q$.
\begin{enumerate}[label=\textup{(\roman*)}]
\item Every apex $\tau_\nu$ is a strict inflection point.
\item If $\sigma_+=-\sigma_-$, then $R$ is an extremal plateau; it is a
strict local extremum exactly when $p=q$.
\item If $\sigma_+=\sigma_-=:\sigma$, then $\Anet$ has no local extremum on
$R$. Its concavity has sign $-\sigma$ immediately to the left and $+\sigma$
immediately to the right. If $p=q$ the point is a strict inflection; if
$p<q$, no point of $R$ is a strict inflection, and the change of concavity
is carried by the whole plateau.
\item Strict local extrema occur exactly at isolated crossing zeros, and
plateaus of extrema exactly at crossing runs of positive length; every point
of such a plateau is a non-strict local extremum.
\end{enumerate}
\end{lemma}

\begin{proof}
On the descending side of the left excursion,
$\sgn\Anet'=\sigma_-$ and $\sgn\Anet''=-\sigma_-$. On the ascending side of
the right excursion, $\sgn\Anet'=\sigma_+$ and $\sgn\Anet''=\sigma_+$. On
$R$, $\Anet'=0$. The assertions follow, and at every apex the sign of
$\Anet''$ reverses.
\end{proof}

\section{Internal correlations of an excursion}
\label{sec:internal_corr}

For a real block $x_1,\dots,x_L$ put
\begin{equation}
\label{eq:block_def}
S_t=\sum_{j=1}^t x_j,\qquad
C_h=\sum_{j=1}^{L-h}x_jx_{j+h}.
\end{equation}

\begin{lemma}[Triangular identity]
\label{lem:triangular-corr}
For every real finite sequence,
\begin{equation}
\label{eq:general-triangular}
\sum_{t=1}^{L-1}S_t^2+\sum_{h=1}^{L-1}hC_h
=S_L\sum_{t=1}^{L-1}S_t.
\end{equation}
\end{lemma}

\begin{proof}
Expanding the partial sums,
\[
\sum_{t=1}^{L}S_t^2
=\sum_{i,j=1}^{L}x_ix_j\bigl(L+1-\max(i,j)\bigr).
\]
Write $\max(i,j)=\tfrac12(i+j+|i-j|)$. The part containing $|i-j|$ is
$\tfrac12\sum_{i,j}x_ix_j|i-j|=\sum_{h=1}^{L-1}hC_h$. The remaining part is
$(L+1)S_L^2-S_L\sum_i ix_i$, and $\sum_i ix_i=LS_L-\sum_{t=1}^{L-1}S_t$.
Hence
$\sum_{t=1}^{L}S_t^2=S_L^2+S_L\sum_{t<L}S_t-\sum_hhC_h$; subtract $S_L^2$.
\end{proof}

\begin{corollary}[Bridge case]
\label{cor:bridge}
If $S_L=0$, then
\begin{equation}
\label{eq:bridge-triangular}
\sum_{t=1}^{L-1}S_t^2=-\sum_{h=1}^{L-1}hC_h.
\end{equation}
\end{corollary}

\begin{remark}[Hermitian form]
\label{rem:hermitian}
For complex $x_j$ and $C_h=\sum_{j=1}^{L-h}x_j\overline{x_{j+h}}$, the same
computation gives
\[
\sum_{t=1}^{L-1}|S_t|^2+\sum_{h=1}^{L-1}h\Re C_h
=\Re\!\left(\overline{S_L}\sum_{t=1}^{L-1}S_t\right).
\]
\end{remark}

For a complete excursion orient the increments by
\begin{equation}
\label{eq:orientation}
x_j=\sigma_\nu\mu(x_{0,\nu}+j),\qquad
S_t=|M(x_{0,\nu}+t)|.
\end{equation}
Then $S_t>0$ for $1\le t<\ell_\nu$ and $S_{\ell_\nu}=0$. Put
\begin{equation}
\label{eq:Q_def}
Q_\nu=\sum_{t=1}^{\ell_\nu-1}S_t^2,\qquad
C_h^{(\nu)}=\sum_{j=1}^{\ell_\nu-h}x_jx_{j+h},
\qquad\text{so that}\qquad
Q_\nu=-\sum_{h=1}^{\ell_\nu-1}hC_h^{(\nu)}.
\end{equation}

\begin{lemma}[Zeroth moment and extreme lag]
\label{lem:zeroth-corr}
\label{lem:extreme-lag}
With $q_\nu=\sum_jx_j^2$,
\begin{equation}
\label{eq:zeroth-corr}
\sum_{h=1}^{\ell_\nu-1}C_h^{(\nu)}=-\frac{q_\nu}{2},\qquad
C_{\ell_\nu-1}^{(\nu)}=-1,
\end{equation}
and $q_\nu$ is even and at least $2$.
\end{lemma}

\begin{proof}
$0=(\sum_jx_j)^2=q_\nu+2\sum_hC_h^{(\nu)}$. The first oriented step is $+1$;
immediately before the first return the path is at height $1$, so the last
step is $-1$. The bridge condition forces equally many positive and negative
active steps.
\end{proof}

\begin{proposition}[Quadratic decomposition by excursions]
\label{prop:quadratic-decomp}
For real $x\ge1$ let $u(x)$ and $L_x$ be as in
Definition~\ref{def:Ms-terminal}, and for $1\le j\le L_x$ put
$y_j=\mu(u(x)+j)$ and $C_h^{\rm term}=\sum_{j=1}^{L_x-h}y_jy_{j+h}$. Then
\begin{equation}
\label{eq:quadratic-decomp-general}
\sum_{n\le x}M(n)^2
=\sum_{\nu:x_{1,\nu}\le u(x)}Q_\nu+Q_{\rm term}(x),
\end{equation}
\begin{equation}
\label{eq:terminal-Q-exact}
Q_{\rm term}(x)=M(\lfloor x\rfloor)^2
+M(\lfloor x\rfloor)\sum_{t=1}^{L_x-1}M(u(x)+t)
-\sum_{h=1}^{L_x-1}hC_h^{\rm term},
\end{equation}
with $Q_{\rm term}(x)=0$ if $L_x=0$. At a zero $x$ of $M$,
\begin{equation}
\label{eq:quadratic-decomp-zero}
\sum_{n\le x}M(n)^2
=\sum_{x_{1,\nu}\le x}Q_\nu
=-\sum_{x_{1,\nu}\le x}\sum_hhC_h^{(\nu)}.
\end{equation}
\end{proposition}

\begin{proof}
Every integer up to $u(x)$ lies either in a zero run or in a unique excursion
interior. Apply Lemma~\ref{lem:triangular-corr} to the terminal block and
add its final term $S_{L_x}^2$.
\end{proof}

\begin{lemma}[Elementary bound for an open terminal block]
\label{lem:terminal-Q-elementary}
Let $u$ be a zero of $M$ and $L\ge0$. Then $|M(u+r)|\le r$ for
$1\le r\le L$, and
\begin{equation}
\label{eq:terminal-Q-elementary}
\sum_{r=1}^{L}M(u+r)^2\le\frac{L(L+1)(2L+1)}6 .
\end{equation}
In particular $Q_{\rm term}(x)\ll L_x^3$. The bridge bound
\eqref{eq:Q_upper}, whose leading coefficient is $1/12$ rather than $1/3$,
applies only to complete excursions.
\end{lemma}

\begin{proof}
$M(u+r)=\sum_{j\le r}\mu(u+j)$ with $|\mu|\le1$; square and sum.
\end{proof}


Thus the quadratic sum sees only lag correlations internal to excursions:
pairs of M\"obius increments in different completed excursions do not
contribute.

Let
\begin{equation}
\label{eq:beta_def}
\beta_\nu=\sum_{t=1}^{\ell_\nu-1}S_t=\sum_{n\in I_\nu}|M(n)|.
\end{equation}
Cauchy--Schwarz gives
\begin{equation}
\label{eq:beta-Q}
\beta_\nu^2\le(\ell_\nu-1)Q_\nu,
\end{equation}
with equality exactly for a flat excursion of height $1$. To separate
support from sign geometry, let $P_h^{(\nu)}=\#\{j:x_jx_{j+h}\ne0\}$ and let
$D_h^{(\nu)}$ count the opposite-sign pairs among them. Then
\begin{equation}
\label{eq:C-P-D}
C_h^{(\nu)}=P_h^{(\nu)}-2D_h^{(\nu)},\qquad
Q_\nu=\sum_{h=1}^{\ell_\nu-1}h\bigl(2D_h^{(\nu)}-P_h^{(\nu)}\bigr).
\end{equation}

\begin{proposition}[Sharp bounds at fixed length and height]
\label{prop:sharp-LH}
Let $S_0=S_L=0$, $S_t\ge1$ for $0<t<L$, $S_t-S_{t-1}\in\{-1,0,1\}$ and
$H=\max S_t$. Then $L\ge2H$ and
\begin{align}
H^2+L-2H&\le\beta\le H(L-H),\label{eq:beta-LH}\\
\frac{2H^3+H}{3}+L-2H&\le Q\le
\frac{2H^3+H}{3}+(L-2H)H^2.\label{eq:Q-LH}
\end{align}
All four bounds are attained. In particular
\begin{equation}
\label{eq:height_Q_min}
Q\ge\frac{2H^3+H}{3},\qquad Q\ge L-1.
\end{equation}
\end{proposition}

\begin{proof}
Before and after the first visit to $H$ the path visits each level
$1,\dots,H-1$. Together with one visit to $H$, these $2H-1$ designated
interior times contribute $2(1+\cdots+(H-1))+H=H^2$ to $\beta$ and
$2\sum_{j<H}j^2+H^2=(2H^3+H)/3$ to $Q$. The remaining $L-2H$ interior times
lie at heights between $1$ and $H$; placing them all at height $1$, or all
at height $H$, attains the bounds.
\end{proof}

\begin{corollary}[Sharp fixed-length bounds]
\label{cor:fixed-L}
For $L\ge2$,
\[
\beta\le\left\lfloor\frac{L^2}{4}\right\rfloor,\qquad
Q\le
\begin{cases}
L(L^2+2)/12,&L\text{ even},\\
L(L^2-1)/12,&L\text{ odd},
\end{cases}
\]
and in particular
\begin{equation}
\label{eq:Q_upper}
Q\le\frac{L^3+2L}{12}.
\end{equation}
Both bounds are attained by the tent path, with one horizontal step at
maximal height when $L$ is odd.
\end{corollary}

\begin{proof}
Use $S_t\le\min(t,L-t)$ and sum this bound and its square.
\end{proof}

\section{Where \texorpdfstring{$\zeta$}{zeta} enters}
\label{sec:zeta}

The black curve depends on $M$ nonlinearly, through $\sqrt{H_\nu}$ and
through the selection of $x_{0,\nu},\tau_\nu,x_{1,\nu}$, and no closed
expression of its Mellin transform in terms of $\zeta$ is known. The nodes
$x_{0,\nu},x_{1,\nu}$ are zeros of $M$. If RH holds and all zeros are
simple, there is a sequence $T_k\to\infty$ such that, for non-integer $x$,
\begin{equation}
\label{eq:explicitM}
  M(x)=\lim_{k\to\infty}\sum_{|\gamma|<T_k}\frac{x^\rho}{\rho\,\zeta'(\rho)}-2
  +\sum_{n\ge1}\frac{(-1)^{n-1}(2\pi/x)^{2n}}{(2n)!\,n\,\zeta(2n+1)}
\end{equation}
\cite[\S14.27]{Titchmarsh1986}. For $\widetilde{A}(x):=\int_0^xM(t)\,dt$ one
has, for $\Re s>1$,
\begin{equation}
\label{eq:mellin_A}
  \int_1^\infty \widetilde{A}(x)x^{-s-2}\,dx=\frac{1}{s(s+1)\zeta(s)}.
\end{equation}
Shifting the contour formally, the residues are
\begin{equation}
\label{eq:explicitA}
  \sum_\rho\frac{x^{\rho+1}}{\rho(\rho+1)\zeta'(\rho)}-2x+12
  +(\text{residues at the trivial zeros}),
\end{equation}
where $x/\zeta(0)=-2x$ is the residue at $s=0$ and
$1/((-1)\zeta(-1))=12$ the residue at $s=-1$. The factor $(\rho+1)^{-1}$
improves convergence, but RH and simplicity alone do not control
$1/\zeta'(\rho)$ well enough to turn \eqref{eq:explicitA} into an identity
with an untruncated sum; results of this type require additional input of
Gonek--Hejhal type \cite{Ng2004}. The area of $\Msk$ differs from
$\widetilde A$ by explicit local terms.

\begin{proposition}[Dipole identity]
\label{prop:dipoles}
For each excursion let its \emph{legs} be the integer sets
$J=(u_\nu,\tau_\nu]\cap\mathbb{Z}$ and
$(\tau_\nu,v_\nu]\cap\mathbb{Z}$, with midpoints $c_J=(u_\nu+\tau_\nu)/2$
and $(\tau_\nu+v_\nu)/2$. Then for every zero $v_k$,
\begin{equation}
\label{eq:dipoles}
  \int_0^{v_k}\Msk(t)\,dt=\widetilde{A}(v_k)
  +\sum_{J\subset(0,v_k]}\ \sum_{m\in J}\mu(m)\,(m-c_J).
\end{equation}
\end{proposition}

\begin{proof}
Since $M$ is constant on each $[n,n+1)$, $\int_a^bM(t)\,dt=\sum_{n=a}^{b-1}M(n)$
for integers $a<b$. On an ascent leg, $\int_{u}^{\tau}\Msk=\sigma H a/2$ and
$\int_u^\tau M=\sum_{m=u+1}^{\tau}\mu(m)(\tau-m)$, while
$\sigma H=\sum_{m=u+1}^{\tau}\mu(m)$; the difference is
$\sum_m\mu(m)(a/2-\tau+m)=\sum_m\mu(m)(m-c_J)$. On a descent leg,
$\int_\tau^v\Msk=\sigma Hb/2$ and
$\int_\tau^vM=b\,\sigma H+\sum_{m=\tau+1}^{v}\mu(m)(v-m)$ with
$\sum_{m=\tau+1}^{v}\mu(m)=-\sigma H$; the difference is
$\sum_m\mu(m)(m-(v-b/2))$, and $v-b/2=(\tau+v)/2=c_J$. On zero runs
both integrands vanish. Sum over legs.
\end{proof}

\section{Bounds for the area curves}
\label{sec:bounds}

\begin{lemma}[Transfer]
\label{lem:transfer}
Let $x\ge1$ and let $x^*\ge x$ be a zero of $M$ (in a computation truncated
at $N\ge x$, take $x^*=N$ and the terminal skeleton of
Definition~\ref{def:Ms-terminal}). If $|M(n)|\le C_0\,n^{1/2}$ for all
$1\le n\le x^*$ (so necessarily $C_0\ge1$), then
\begin{equation}
\label{eq:transfer_bound}
  |\Ms(t)|\le C_0^{1/2}t^{1/4}\quad(0\le t\le x),\qquad
  \Ap(x)+|\Am(x)|\le\tfrac45\,C_0^{1/2}x^{5/4}.
\end{equation}
\end{lemma}

\begin{proof}
For $t\ge\tau_\nu$, $|\Ms(t)|\le\sqrt{H_\nu}\le C_0^{1/2}\tau_\nu^{1/4}\le
C_0^{1/2}t^{1/4}$. For $u_\nu\le t\le\tau_\nu$ the chord from $(u_\nu,0)$ to
$(\tau_\nu,\sqrt{H_\nu})$ lies below the chord of $C_0^{1/2}t^{1/4}$ over
$[u_\nu,\tau_\nu]$, whose endpoints $C_0^{1/2}u_\nu^{1/4}\ge0$ and
$C_0^{1/2}\tau_\nu^{1/4}\ge\sqrt{H_\nu}$ (as $\tau_\nu\le x^*$) are at
least as high; that chord lies below $C_0^{1/2}t^{1/4}$ by concavity.
Integrate.
\end{proof}

\begin{lemma}[A concave majorant]
\label{lem:concave-majorant}
Let $\phi(y)=(\log y)^{3/5}(\log\log y)^{-1/5}$ and, for $\kappa>0$,
$F(y)=y^{1/2}e^{-\kappa\phi(y)}$. There is $y_0=y_0(\kappa)$ such that $F$
is increasing and concave on $[y_0,\infty)$.
\end{lemma}

\begin{proof}
With $y=e^s$ and $g(s)=s^{3/5}(\log s)^{-1/5}$, one has
$F=\exp(s/2-\kappa g(s))$. Put $h(s)=\frac12-\kappa g'(s)$. Then
\[
yF'(y)=F\,h(s),\qquad
y^2F''(y)=F\bigl(h(s)^2-h(s)-\kappa g''(s)\bigr).
\]
Since $g'(s)\to0$ and $g''(s)\to0$, $h(s)\to\frac12$ and
$h^2-h-\kappa g''\to-\frac14$. Hence $F'>0$ and $F''<0$ for large $y$.
\end{proof}

\begin{corollary}
\label{cor:uncond}
From the Vinogradov--Korobov bound \cite{Walfisz1963}
\begin{equation}
\label{eq:VK_M}
  M(x)\ll x\exp\bigl(-c\,\phi(x)\bigr)
\end{equation}
one obtains
\begin{equation}
\label{eq:uncond_area}
  \Ap(x)+|\Am(x)|\ll x^{3/2}\exp\bigl(-\tfrac c2\,\phi(x)\bigr).
\end{equation}
\end{corollary}

\begin{proof}
Let $K$ be such that $|M(y)|\le Ky\,e^{-c\phi(y)}$ for $y\ge y_1$, and let
$F$ be as in Lemma~\ref{lem:concave-majorant} with $\kappa=c/2$, so that
$\sqrt{|M(y)|}\le K^{1/2}F(y)$ for $y\ge y_1$. Only finitely many excursions
start below $Y_0:=\max(y_0,y_1)$; they are finite, so they contribute
$O(1)$ to the area. Consider an excursion with $u_\nu\ge Y_0$. Its apex
satisfies $\sqrt{H_\nu}=\sqrt{|M(\tau_\nu)|}\le K^{1/2}F(\tau_\nu)$; the
bound is taken at the apex $\tau_\nu$, which lies in the excursion, and not
at its future endpoint $v_\nu$. For $t\ge\tau_\nu$ monotonicity of $F$ gives
$|\Ms(t)|\le\sqrt{H_\nu}\le K^{1/2}F(t)$. For $u_\nu\le t\le\tau_\nu$,
$|\Ms(t)|$ is the chord from $(u_\nu,0)$ to $(\tau_\nu,\sqrt{H_\nu})$. This
chord lies below the chord from $(u_\nu,K^{1/2}F(u_\nu))$ to
$(\tau_\nu,K^{1/2}F(\tau_\nu))$, which has both endpoints at least as high,
and the latter lies below $K^{1/2}F$ on $[u_\nu,\tau_\nu]\subset[y_0,\infty)$
by concavity. Hence $|\Ms(t)|\le K^{1/2}F(t)$ for all $t\ge Y_0$ outside the
finitely many early excursions, and since $F$ is increasing,
\[
\Ap(x)+|\Am(x)|\le O(1)+K^{1/2}\int_{Y_0}^xF(t)\,dt
\le O(1)+K^{1/2}xF(x)=O(1)+K^{1/2}x^{3/2}e^{-\frac c2\phi(x)}.
\qedhere
\]
\end{proof}

\begin{corollary}
\label{cor:rh}
If RH holds, then
\begin{equation}
\label{eq:rh_area}
  \Ap(x)+|\Am(x)|\ll_\eps x^{5/4+\eps}.
\end{equation}
Consequently, if $\Ap(x)+|\Am(x)|\gg x^{5/4+\delta}$ along an unbounded
sequence for some fixed $\delta>0$, then RH is false.
\end{corollary}

\begin{remark}[Endpoint-indexed bounds]
\label{rem:onedirection}
From $\ell_\nu\ge2H_\nu$ and $\sqrt{H_\nu}\ell_\nu\ll v_\nu^{5/4+\eps}$ one
deduces only $H_\nu\ll v_\nu^{5/6+O(\eps)}$. This bound is indexed by the
\emph{future} endpoint $v_\nu$ of the excursion; it gives no bound for $M(x)$
in terms of $x$ at interior points, and hence no zero-free half-plane.
Whether \eqref{eq:rh_area} is equivalent to RH is open.
\end{remark}

\begin{remark}[Oscillation of the heights]
\label{rem:omega}
The Mertens conjecture is false \cite{OdlyzkoTeRiele1985,KotnikTeRiele2006};
the current bounds are due to Hurst~\cite{Hurst18}, who proves
$\limsup M(x)/\sqrt{x}>1.826054$ and
$\liminf M(x)/\sqrt{x}<-1.837625$. The excursion containing a point $n$ with
$M(n)>1.826054\sqrt n$ has $u_\nu<n$, so $H_\nu>1.826054\,u_\nu^{1/2}$ for
infinitely many positive excursions, and similarly with $1.837625$ for
negative ones. A lower bound indexed by $v_\nu$ does not follow, since
$v_\nu$ may be much larger than $n$.
\end{remark}

\section{Global criteria equivalent to RH}
\label{sec:criteria}

The Mellin identity and the global $L^1/L^2$ equivalences below are
consequences of Littlewood's criterion and partial summation; the excursion
decompositions and pruning lemmas repackage them in terms of excursions.
For $\Re s>1$,
\begin{equation}
\label{eq:mellin_M}
  \int_1^\infty M(x)x^{-s-1}\,dx=\frac{1}{s\,\zeta(s)} .
\end{equation}
If $\int_1^\infty|M(x)|x^{-\varsigma-1}dx<\infty$ for every $\varsigma>1/2$, the
left side continues analytically to $\Re s>1/2$, so $\zeta$ has no zeros
there.

\begin{proposition}[$L^1$ criterion]
\label{prop:L1}
RH holds if and only if, for every $\eps>0$,
\begin{equation}
\label{eq:L1_criterion}
  \sum_{v_\nu\le x}\beta_\nu+\beta_{\rm term}(x)=\sum_{n\le x}|M(n)|
  \ll x^{3/2+\eps}.
\end{equation}
\end{proposition}

\begin{proof}
RH gives $M(x)\ll_\eps x^{1/2+\eps}$. Conversely, for $\varsigma>1/2$ choose
$0<\eps<\varsigma-1/2$; partial summation gives
$\sum_n|M(n)|n^{-\varsigma-1}<\infty$, and since $M$ is constant on $[n,n+1)$
the remark after \eqref{eq:mellin_M} applies.
\end{proof}

\begin{lemma}[Pruning]
\label{lem:pruning}
$\beta_\nu\le H_\nu(\ell_\nu-1)\le\ell_\nu^2/2$, hence
\begin{equation}
\label{eq:pruning}
  \sum_{v_\nu\le x,\ \ell_\nu\le L_0}\beta_\nu\le\frac{xL_0}{2}.
\end{equation}
\end{lemma}

\begin{corollary}[Reduction to long excursions]
\label{cor:long}
For fixed $\delta\ge0$ and $L_0=x^{1/2-\delta}$, RH holds if and only if
\begin{equation}
\label{eq:long_reduction}
  \sum_{\substack{v_\nu\le x\\ \ell_\nu>L_0}}\beta_\nu
  +\mathbf 1_{\{L_x>L_0\}}\beta_{\rm term}(x)\ll x^{3/2+\eps}
  \quad\text{for every }\eps>0.
\end{equation}
At most $x/L_0$ long excursions are involved.
\end{corollary}

\begin{proof}
If RH holds, the left side is at most $\sum_{n\le x}|M(n)|\ll x^{3/2+\eps}$.
Conversely, the complete excursions with $\ell_\nu\le L_0$ contribute at
most $xL_0/2\le x^{3/2}$ by Lemma~\ref{lem:pruning}, and if $L_x\le L_0$
the terminal block contributes
$\beta_{\rm term}(x)\le\sum_{r\le L_x}r=L_x(L_x+1)/2\le L_0^2\le x$.
Hence $\sum_{n\le x}|M(n)|\ll x^{3/2+\eps}$, and
Proposition~\ref{prop:L1} applies.
\end{proof}

\begin{remark}[The $\sqrt{x}$ frontier]
\label{rem:sqrt_frontier}
The cut may depend on $\eps$: with $L_0=x^{1/2+\eps/2}$ the discarded part is
$\le x^{3/2+\eps/2}/2$. The frontier of pruning is therefore $x^{1/2+o(1)}$;
what requires new information is a cut $x^{\theta}$ with a fixed
$\theta>1/2$. For $\ell_\nu\asymp x$ the trivial bound gives
$\beta_\nu\ll x^2$ against a budget $x^{3/2+\eps}$, and there are only
$O(1)$ such excursions.
\end{remark}

\begin{heuristic}[Renewal tail]
\label{heu:renewal}
If excursion lengths behave like those of a recurrent random walk,
\begin{equation}
\label{eq:renewal_tail}
  \#\{v_\nu\le x:\ \ell_\nu>L\}\asymp\sqrt{x}\,L^{-1/2},
\end{equation}
and with $\beta_\nu\asymp\ell_\nu^{3/2}$ the dyadic range $[L,2L]$
contributes $\asymp\sqrt{x}\,L$ to \eqref{eq:L1_criterion}, so the terminal
scales $L\asymp x$ dominate the $L^1$ sum.
\end{heuristic}

\begin{theorem}[Correlation criterion]
\label{thm:corr-criterion}
RH holds if and only if, for every $\eps>0$,
\begin{equation}
\label{eq:RH-corr}
  \sum_{n\le x}M(n)^2=\sum_{v_\nu\le x}Q_\nu+Q_{\rm term}(x)
  =-\sum_{v_\nu\le x}\sum_{h}hC^{(\nu)}_h+Q_{\rm term}(x)\ll x^{2+\eps}.
\end{equation}
\end{theorem}

\begin{proof}
Under RH, $M(n)^2\ll_\eps n^{1+\eps}$. Conversely, by Cauchy--Schwarz,
\[
\sum_{n\le x}|M(n)|\le x^{1/2}\Bigl(\sum_{n\le x}M(n)^2\Bigr)^{1/2}\ll x^{3/2+\eps/2};
\]
apply
Proposition~\ref{prop:L1}.
\end{proof}

\begin{corollary}[Local sufficient conditions]
\label{cor:local-corr-sufficient}
\label{cor:pruned_local}
Suppose that for every $\eps>0$ there is $C_\eps$ such that
\begin{equation}
\label{eq:pruned_Q}
  Q_\nu\le C_\eps\,\ell_\nu^{2+\eps}\quad\text{for every complete excursion with }
  \ell_\nu>v_\nu^{1/2},
\end{equation}
Then RH holds.
\end{corollary}

\begin{proof}
Split the completed excursions at $\ell_\nu=x^{1/2}$. If
$\ell_\nu\le x^{1/2}$, then \eqref{eq:Q_upper} gives
$Q_\nu\le\ell_\nu x/12+\ell_\nu/6$, and disjointness bounds the total by
$x^2/12+x/6$. If $\ell_\nu>x^{1/2}$ then $\ell_\nu>v_\nu^{1/2}$ since
$v_\nu\le x$, and the hypothesis bounds the total by
$C_\eps\sum\ell_\nu^{2+\eps}\le C_\eps(\sum\ell_\nu)^{2+\eps}\le C_\eps x^{2+\eps}$.
The terminal block is not a complete excursion, and \eqref{eq:Q_upper} does
not apply to it. By Lemma~\ref{lem:terminal-Q-elementary},
$Q_{\rm term}(x)\le L_x(L_x+1)(2L_x+1)/6$, which is $O(x^{3/2})$ if $L_x\le x^{1/2}$. If $L_x>x^{1/2}$, let $\nu$ be the
complete excursion containing $x$, starting at $u(x)$; then
$Q_{\rm term}(x)\le Q_\nu$. Fix $A=1+\eps$. By
Corollary~\ref{cor:KP-local-returns}, $v_\nu\le u(x)^A\le x^A$ once $x$ is
large, so $\ell_\nu\le x^A$. If $\ell_\nu>v_\nu^{1/2}$ the hypothesis gives
$Q_\nu\le C_\eps x^{A(2+\eps)}$; otherwise \eqref{eq:Q_upper} gives
$Q_\nu\ll\ell_\nu^3\le v_\nu^{3/2}\le x^{3A/2}$. In all cases
$Q_{\rm term}(x)\ll_\eps x^{2+4\eps}$. This step uses the arithmetic
input of Corollary~\ref{cor:KP-local-returns}: for a general family of
disjoint intervals, the future return of the terminal block is not
controlled by $x$, and a bound for complete excursions would not control
$Q_{\rm term}(x)$. Apply
Theorem~\ref{thm:corr-criterion} with $\eps$ replaced by $4\eps$.
\end{proof}

\begin{heuristic}[Ballistic excursions]
\label{heu:local_fails}
In a model where, after a zero of $M$, the next values of $\mu$ behave
locally at random with squarefree density $6/\pi^2$, the pattern ``$k$
non-zero steps of one sign followed by $k$ of the other'' follows a given
zero with probability $\ge c_0^{-k}$. If excursions with $H_\nu=k$ and
$\ell_\nu\asymp k$ occur for unbounded $k$, then by \eqref{eq:height_Q_min}
$Q_\nu/\ell_\nu^{2+\eps}\gg k^{1-\eps}\to\infty$, so the unpruned version of
\eqref{eq:pruned_Q} fails for every $\eps\in(0,1)$. Such excursions have
$\ell_\nu\ll v_\nu^{1/2}$ and do not affect the pruned condition
\eqref{eq:pruned_Q}; this is why the pruning in
Corollary~\ref{cor:pruned_local} is needed.
\end{heuristic}

\begin{theorem}[Energy criterion]
\label{thm:energy}
With $u$ the last zero $\le x$ and $H_{\rm run}(x)=\max_{u<n\le x}|M(n)|$,
define
\begin{equation}
\label{eq:energy_def}
  \En(x):=\frac14\Bigl(\sum_{v_\nu\le x}H_\nu\ell_\nu+(x-u)H_{\rm run}(x)\Bigr).
\end{equation}
For a complete excursion let $J$ be the chord joining
$(u_\nu,\Anet(u_\nu))$ to $(v_\nu,\Anet(v_\nu))$. By \eqref{eq:jump},
$J'=\sigma_\nu\sqrt{H_\nu}/2$, so $\tfrac14H_\nu\ell_\nu=\int_{I_\nu}(J')^2$
is the Dirichlet energy of the chord. Then
\begin{equation}
\label{eq:energy_sandwich}
  \sum_{n\le x}|M(n)|\le4\En(x)\le x\max_{n\le x}|M(n)|,
\end{equation}
and RH holds if and only if $\En(x)\ll x^{3/2+\eps}$ for every $\eps>0$.
\end{theorem}

\begin{proof}
$\beta_\nu\le H_\nu\ell_\nu$, $\beta_{\rm term}(x)\le(x-u)H_{\rm run}(x)$ and
$\sum\ell_\nu+(x-u)\le x$; apply Proposition~\ref{prop:L1}.
\end{proof}

\section{Lacunary excursion families and horizon-localized criteria}
\label{sec:lacunary}

\subsection{Geometric counts}
\label{subsec:geometric}

The lemmas of this subsection use only that excursions are disjoint integer
intervals ordered by their starting points, the next one starting at or
after the end of the previous one. They therefore hold verbatim for every
sequence of disjoint integer intervals, deterministic or random, and in
particular pathwise for any stochastic model of excursions
(Section~\ref{subsec:extreme}). For a family $\F$ of excursions and $X\ge2$
put
\[
E_\F(X):=\#\{\nu\in\F:x_{0,\nu}\le X\}.
\]
For a function $\mathfrak h:[0,\infty)\to[0,\infty]$, the family of
excursions \emph{crossing the horizon} $h$ is
\[
\F_{\mathfrak h}:=\{\nu:x_{1,\nu}>\mathfrak h(x_{0,\nu})\}.
\]
For $A>1$ the excursion $I_\nu$ is \emph{$A$-exceptional} if
$x_{1,\nu}>x_{0,\nu}^A$, i.e.\ $\nu\in\F_{x^A}$, and \emph{$A$-ordinary}
otherwise; for $A=2$ we say \emph{square-exceptional} and
\emph{square-ordinary}. For $\lambda>1$ it is \emph{$\lambda$-long} if
$x_{1,\nu}>\lambda x_{0,\nu}$. We write $E_A=E_{\F_{x^A}}$,
$E_\square=E_2$ and $E_\lambda=E_{\F_{\lambda x}}$; the context separates
$E_A$ ($A$ an exponent) from $E_\lambda$ ($\lambda$ a factor).

\begin{lemma}[Power lacunarity]
\label{lem:power-lacunary}
For every $A>1$ and $X\ge2$,
\begin{equation}
\label{eq:power-lacunary-count}
E_A(X)\le2+\frac1{\log A}\log\frac{\log X}{\log2}.
\end{equation}
In particular
\begin{equation}
\label{eq:square-exception-count}
E_\square(X)\le2+\log_2\frac{\log X}{\log2},\qquad
E_\square(X)\le4\log\log X\quad(X\ge e^e).
\end{equation}
\end{lemma}

\begin{proof}
No excursion starts at $1$, since $M(1)=1$. List the $A$-exceptional starts
$\ge2$ as $2\le a_1<a_2<\cdots$. The next excursion starts at or after the
end of the preceding one, so $a_{j+1}\ge x_{1,j}>a_j^A$ and
$\log a_j\ge A^{j-1}\log2$. If $a_j\le X$ then
$j\le1+\log(\log X/\log2)/\log A$; add $1$ for the excursion $(0,2]$. For the
second bound in \eqref{eq:square-exception-count},
$\log_2(\log X/\log2)=(\log\log X-\log\log2)/\log2\le1.443\log\log X+0.529$.
\end{proof}

The next lemma is the partition of the square horizon
$[x_0,x_0^2]=[x_0,cx_0^2]\cup[cx_0^2,x_0^2]$: the upper part is accessible
only to a doubly lacunary family, for every fixed $c>0$.

\begin{lemma}[Fractional square horizon]
\label{lem:partition}
For $0<c\le1$ and $X\ge2$,
\begin{equation}
\label{eq:partition-count}
E_{\F_{cx^2}}(X)=\#\{\nu:x_{0,\nu}\le X,\ x_{1,\nu}>cx_{0,\nu}^2\}
\le\frac2c+2+\log_2\frac{\log X}{\log2}.
\end{equation}
Consequently, the range $cx_0^2<t\le x_0^2$ of the excursions starting at
most $X$ is visited by at most $O_c(\log\log X)$ excursions.
\end{lemma}

\begin{proof}
At most $\lceil2/c\rceil\le2/c+1$ members start below $2/c$. List the others
as $a_1<a_2<\cdots$ and put $b_j=ca_j\ge2$. Then
$a_{j+1}\ge x_{1,j}>ca_j^2$, so
\[
b_{j+1}=ca_{j+1}>c^2a_j^2=b_j^2,
\]
and $\log b_j\ge2^{j-1}\log2$. Since $b_j\le cX\le X$, the number of such
$j$ is at most $1+\log_2(\log X/\log2)$. An excursion containing a point
$t>cx_0^2$ has $x_1\ge t>cx_0^2$.
\end{proof}

\begin{lemma}[Superlinear horizons]
\label{lem:superlinear}
Let $\psi:[1,\infty)\to(0,\infty)$ satisfy $\psi(x)\to\infty$. Then
\[
E_{\F_{x\psi(x)}}(X)
=\#\{\nu:x_{0,\nu}\le X,\ x_{1,\nu}>x_{0,\nu}\psi(x_{0,\nu})\}
=o(\log X).
\]
\end{lemma}

\begin{proof}
Given $K>0$ choose $x_K$ with $\psi(x)\ge e^K$ for $x\ge x_K$. At most
$x_K+1$ members start below $x_K$. For consecutive members
$a_j<a_{j+1}$ starting at or above $x_K$,
$a_{j+1}\ge x_{1,j}>a_j\psi(a_j)\ge e^Ka_j$, so $\log a_j$ increases by more
than $K$ at each step and there are at most $1+\log X/K$ of them. Hence
$\limsup E(X)/\log X\le1/K$ for every $K$.
\end{proof}

\begin{lemma}[Linear lacunarity]
\label{lem:linear-lacunary}
For every $\lambda>1$ and $X\ge2$,
\begin{equation}
\label{eq:linear-lacunary-count}
E_\lambda(X)=\#\{\nu:x_{0,\nu}\le X,\ x_{1,\nu}>\lambda x_{0,\nu}\}
\le2+\frac{\log(X/2)}{\log\lambda}.
\end{equation}
\end{lemma}

\begin{proof}
For starts $a_j\ge2$, $a_{j+1}\ge x_{1,j}>\lambda a_j$, hence
$a_j>2\lambda^{j-1}$ for $j\ge2$; add $1$ for $(0,2]$.
\end{proof}

All four bounds are sharp up to constants for general families of disjoint
intervals, as geometric and doubly geometric chains of intervals show.

\subsection{Arithmetic input: Kaczorowski--Pintz}

We separate the published statement, its translation from $M_0=M+2$ to
$M$, and the consequences for excursions.

\begin{theorem}[Kaczorowski--Pintz {\cite{KP86}}]
\label{thm:KP-quoted}
Let $M_0(x)=M(x)-1/\zeta(0)=M(x)+2$, the function $M_a$ of \cite{KP86} at
$a=0$. Let $V(f,Y)$ be defined as in \cite{KP86}, through the largest number of
points of $[1,Y]$ at which $f$ is nonzero with alternating signs; the number
of sign changes of $f$ on $[1,Y]$ is at least $V(f,Y)-1$, which does not
affect the asymptotic statements below. Then:
\begin{enumerate}[label=\textup{(\roman*)}]
\item \cite[Corollary~4, p.~178, (2.13)--(2.15)]{KP86}
\begin{equation}
\label{eq:KP-signchanges}
\liminf_{Y\to\infty}\frac{V(M_0,Y)}{\log Y}\ge\frac{\gamma_1}{\pi},
\end{equation}
with combined oscillation of size $x^{\Theta-\eta}$ for every $\eta>0$, in
the sense of their definition following (1.1): the alternating points $t_j$
counted by $V$ can be chosen with $|M_0(t_j)|\ge t_j^{\Theta-\eta}$;
\item Corollary~4 is deduced from their Corollary~1, which asserts relations
(2.1)--(2.3) of their Theorem~1 \cite[pp.~175--176]{KP86}; in particular,
for every $\delta>0$ and $Y>Y_0(\delta)$, the interval $[Y^{1-\delta},Y]$
contains two points $t_+,t_-$ with
$M_0(t_+)\ge t_+^{\Theta-\eta}$ and $M_0(t_-)\le-t_-^{\Theta-\eta}$.
\end{enumerate}
\end{theorem}

\begin{lemma}[From $M_0$ to $M$]
\label{lem:KP-translation}
If $|M_0(t)|>2$, then $\sgn M(t)=\sgn M_0(t)$ and
$|M(t)|\ge|M_0(t)|-2$. In particular, for points with
$|M_0(t)|\ge t^{\Theta-\eta}$ and $t$ large,
$\sgn M(t)=\sgn M_0(t)$ and $|M(t)|\ge\frac12t^{\Theta-\eta}$.
\end{lemma}

\begin{proof}
$M_0>2$ gives $M>0$; $M_0<-2$ gives $M<-4$. The second assertion uses
$t^{\Theta-\eta}\to\infty$ for $\eta<\Theta$.
\end{proof}

\begin{corollary}[Distinct large excursions]
\label{cor:KP-large-excursions}
\label{thm:KP-M}
If $\Theta>1/2$ and $0<\eta<\Theta-1/2$, then
\begin{equation}
\label{eq:KP-large-excursions}
N_\eta(Y)\ge\left(\frac{\gamma_1}{\pi}-o(1)\right)\log Y,
\end{equation}
where $N_\eta(Y)$ counts distinct excursions starting at most $Y$ that
contain a point $t\le Y$ with $|M(t)|\ge\tfrac12 t^{\Theta-\eta}$.
\end{corollary}

\begin{proof}
Take the alternating points of Theorem~\ref{thm:KP-quoted}(i) in $[1,Y]$,
discarding the boundedly many below the threshold of
Lemma~\ref{lem:KP-translation}. By that lemma, consecutive points $t<t'$
have $M(t)$ and $M(t')$ of opposite signs. Since $M$ is integer valued with
steps $0,\pm1$, it vanishes at some integer in $(t,t')$, so $t$ and $t'$ lie
in different excursions. Each point with $M(t)\ne0$ belongs to exactly one
excursion, and excursions are ordered, so the excursions containing the
ordered points are pairwise distinct. Their starts are below the points,
hence at most $Y$, and $|M(t)|\ge\frac12t^{\Theta-\eta}$ at each point.
\end{proof}

\begin{corollary}[Multiplicative localization of Mertens returns]
\label{cor:KP-local-returns}
\label{cor:finiteness_expanding}
For every fixed $A>1$ there exists $U_A$ such that every excursion starting
at $u\ge U_A$ ends at $v\le u^A$. Equivalently, $M$ has only finitely many
$A$-exceptional excursions; in particular only finitely many
square-exceptional ones.
\end{corollary}

\begin{proof}
Put $\delta=1-1/A$ and $Y=u^A$, so that $[Y^{1-\delta},Y]=[u,u^A]$, and fix
$0<\eta<\Theta$ (possible since $\Theta\ge1/2$). For $u$ large,
Theorem~\ref{thm:KP-quoted}(ii) and Lemma~\ref{lem:KP-translation} give
$t_+,t_-\in[u,u^A]$ with $M(t_+)>0>M(t_-)$. Neither point lies in $[u,u+1)$,
because $M(u)=0$; so $u<\lfloor t_\pm\rfloor$. Since $M$ changes by
$0,\pm1$ at each integer, it vanishes at some integer $n$ strictly between
$\lfloor t_+\rfloor$ and $\lfloor t_-\rfloor$, and $u<n<u^A$. If
$u=x_{0,\nu}$, the first return satisfies $x_{1,\nu}\le n<u^A$.
\end{proof}

\begin{remark}
A bare sign change of $M_0=M+2$ would not force a zero of $M$. The
decisive point is that the local sign change is combined with oscillation of
size $t^{\Theta-\eta}\to\infty$, so that Lemma~\ref{lem:KP-translation}
applies to both witnesses.
\end{remark}

\begin{remark}[Geometric counts and arithmetic localization]
\label{rem:geom-vs-arith}
For the actual Mertens excursions Corollary~\ref{cor:KP-local-returns} is
stronger than Lemma~\ref{lem:power-lacunary}: each $A$-exceptional family is
finite. The geometric lemmas remain separate results because they are
elementary, independent of $\zeta$, and valid for every ordered family of
disjoint intervals. This is exactly what the localization theorem below and
the extreme-value theorems of Section~\ref{subsec:extreme} use: the former
needs only a count below $(\gamma_1/\pi)\log X$, the latter needs a count
valid pathwise in a model where no analogue of Kaczorowski--Pintz is
available.
\end{remark}

\subsection{Horizon localization}
\label{subsec:horizon}

\begin{theorem}[Horizon-localized RH criterion]
\label{thm:horizon}
\label{thm:sparse-RH}
Let $\mathfrak h:[0,\infty)\to[0,\infty]$ be a function, let $\F$ be a family
of excursions, and suppose that for some $\kappa<\gamma_1/\pi$
\begin{equation}
\label{eq:horizon-hyp}
E_{\F\cup\F_{\mathfrak h}}(X)\le(\kappa+o(1))\log X\qquad(X\to\infty).
\end{equation}
Then RH is equivalent to the following statement: for every $\eps>0$ there
is $C_\eps$ such that
\begin{equation}
\label{eq:restricted-RH}
|M(t)|\le C_\eps t^{1/2+\eps}
\qquad\text{for all }t\in I_\nu,\ \nu\notin\F,\ t\le \mathfrak h(x_{0,\nu}).
\end{equation}
No condition is imposed on the excursions in $\F$, nor on the part of any
excursion lying beyond its horizon $\mathfrak h(x_0)$.
\end{theorem}

\begin{proof}
RH implies \eqref{eq:restricted-RH} by Littlewood's criterion. Conversely,
suppose $\Theta>1/2$ and fix $\eta,\eps>0$ with
$1/2+\eps<\Theta-\eta$. By \eqref{eq:KP-large-excursions} there are at least
$(\gamma_1/\pi-o(1))\log X$ distinct excursions starting at most $X$, each
containing a witness $t$ with $|M(t)|\ge\frac12t^{\Theta-\eta}$. Consider
such an excursion $\nu$ with witness $t$. If $\nu\notin\F$ and
$t\le \mathfrak h(x_{0,\nu})$, the witness lies in the range of
\eqref{eq:restricted-RH}. Otherwise $\nu\in\F$, or $t>\mathfrak h(x_{0,\nu})$, in
which case $x_{1,\nu}>t>\mathfrak h(x_{0,\nu})$ (as $M(t)\ne0$) and $\nu\in\F_{\mathfrak h}$. By
\eqref{eq:horizon-hyp} at least
\[
\Bigl(\frac{\gamma_1}{\pi}-\kappa-o(1)\Bigr)\log X\longrightarrow\infty
\]
witnesses lie in the range of \eqref{eq:restricted-RH}. They belong to
distinct excursions, so they are unbounded, and at them
$t^{\Theta-\eta}/t^{1/2+\eps}\to\infty$, contradicting
\eqref{eq:restricted-RH}.
\end{proof}

The case $\mathfrak h(x)=\infty$ excludes a sparse family $\F$ of whole excursions;
the case $\F=\emptyset$ truncates every excursion at its horizon.

\begin{corollary}[Square horizon]
\label{cor:power-ordinary-RH}
\label{cor:square-truncated-RH}
For every fixed $A>1$, RH is equivalent to \eqref{eq:restricted-RH}
restricted to $A$-ordinary excursions. In particular
\begin{equation}
\label{eq:square-horizon-RH}
\boxed{\ \mathrm{RH}\iff
|M(t)|\ll_\eps t^{1/2+\eps}
\ \text{for all $t$ in excursions with }x_1\le x_0^2,\ }
\end{equation}
and, equivalently, RH holds if and only if
$|M(t)|\ll_\eps t^{1/2+\eps}$ for $x_{0,\nu}<t\le\min\{x_{1,\nu},x_{0,\nu}^2\}$
over all excursions.
\end{corollary}

\begin{proof}
Take $\F=\F_{x^A}$, $\mathfrak h=\infty$, respectively $\F=\emptyset$, $\mathfrak h(x)=x^2$, in
Theorem~\ref{thm:horizon}; Lemma~\ref{lem:power-lacunary} gives
$\kappa=0$.
\end{proof}

The count $E_\square(X)=O(\log\log X)$ is thus what makes the square
horizon admissible. The square horizon can be lowered.

\begin{corollary}[Descending horizons]
\label{cor:descending}
\label{cor:linear-ordinary-RH}
RH is equivalent to Littlewood's bound
$|M(t)|\ll_\eps t^{1/2+\eps}$ on the ranges
\[
x_{0,\nu}<t\le\min\{x_{1,\nu},\mathfrak h(x_{0,\nu})\}
\]
over all excursions, for each of the following horizons:
\begin{enumerate}[label=\textup{(\roman*)}]
\item $\mathfrak h(x)=cx^2$, for any fixed $c>0$;
\item $\mathfrak h(x)=x\psi(x)$, for any $\psi$ with $\psi(x)\to\infty$;
\item $\mathfrak h(x)=\lambda x$, for any fixed $\lambda>\lambda_*:=e^{\pi/\gamma_1}=1.24889\ldots$
\end{enumerate}
The same holds with ``truncate at $\mathfrak h(x_0)$'' replaced by ``discard the
excursions with $x_1>\mathfrak h(x_0)$''. In particular it is enough to control the
excursions with $x_1\le\frac54x_0$.
\end{corollary}

\begin{proof}
Apply Theorem~\ref{thm:horizon} with $\F=\emptyset$ (respectively
$\F=\F_{\mathfrak h}$, $\mathfrak h=\infty$). For (i), $c\ge1$ reduces to
Lemma~\ref{lem:power-lacunary} and $0<c<1$ to Lemma~\ref{lem:partition};
$\kappa=0$. For (ii), Lemma~\ref{lem:superlinear} gives $\kappa=0$. For
(iii), Lemma~\ref{lem:linear-lacunary} gives $\kappa=1/\log\lambda$, and
$1/\log\lambda<\gamma_1/\pi$ exactly when $\lambda>\lambda_*$. Finally
$5/4>\lambda_*$.
\end{proof}

Case (iii) implies (i) and (ii): in both, $\mathfrak h(x_0)\ge\lambda x_0$ for large
$x_0$, so the bound on the larger ranges implies it on the ranges of (iii),
the finitely many remaining excursions affecting only $C_\eps$. The horizon has
been lowered from $x_0^2$ to $\lambda x_0$: the possible failure of
Littlewood's bound, if any, is visible within a bounded multiplicative
distance of the start of an excursion. Quantitatively:

\begin{corollary}[Two-sided localization of large values]
\label{cor:two-sided}
Suppose $\Theta>1/2$, and fix $0<\eta<\Theta-\frac12$ and $\lambda>\lambda_*$.
Among the excursions starting at $x_0\le X$, at least
\[
\Bigl(\frac{\gamma_1}{\pi}-\frac1{\log\lambda}-o(1)\Bigr)\log X
\]
contain a point $t_*$ with $|M(t_*)|\ge\frac12t_*^{\Theta-\eta}$ and
\[
x_0+\tfrac12x_0^{\Theta-\eta}\le t_*\le\lambda x_0 .
\]
\end{corollary}

\begin{proof}
Of the excursions of Corollary~\ref{cor:KP-large-excursions}, those whose
witness exceeds $\lambda x_0$ are $\lambda$-long, hence at most
$\log X/\log\lambda+2$ by Lemma~\ref{lem:linear-lacunary}. For the others,
$M(x_0)=0$ and unit steps give $|M(t_*)|\le t_*-x_0$, so
$t_*-x_0\ge\frac12t_*^{\Theta-\eta}\ge\frac12x_0^{\Theta-\eta}$.
\end{proof}


\begin{remark}[Dyadic erosion and the uniform transition]
\label{rem:dyadic-erosion}
For every fixed integer $j\ge0$, Lemma~\ref{lem:partition} with
$c=2^{-j}$ gives
\[
\#\left\{\nu:x_{0,\nu}\le X,\ x_{1,\nu}>
\frac{x_{0,\nu}^2}{2^j}\right\}=O_{2^j}(\log\log X),
\]
so Theorem~\ref{thm:horizon} successively permits the horizons
\[
x_0^2,\quad \frac{x_0^2}{2},\quad \frac{x_0^2}{4},\quad
\frac{x_0^2}{8},\ldots.
\]
The constant in the fixed-$c$ estimate depends on $c$; hence one may not set
$j=j(x_0)\sim\log_2x_0$ by that lemma alone.  If instead
$2^{j(x)}=x/\psi(x)$ with $\psi(x)\to\infty$, then the corresponding
horizon is $x\psi(x)$ and Lemma~\ref{lem:superlinear} supplies the uniform
$o(\log X)$ crossing count.  Thus the rigorous erosion chain is
\[
x_0^2\to\frac{x_0^2}{2}\to\frac{x_0^2}{4}\to\cdots
\to x_0\psi(x_0)\to\lambda x_0,\qquad\lambda>\lambda_*.
\]
\end{remark}

\subsection{Fixed-mask Catalan compression inside fractional square horizons}
\label{subsec:catalan-square}

The fractional-square localization has a useful interaction with the exact
combinatorics of a completed excursion.  The relevant Catalan time is not
physical time: zero M\"obius increments must first be removed.  This
support-time distinction is essential below.

\begin{definition}[Support-time path and fixed-mask Catalan law]
\label{def:support-catalan}
Let $I=(x_0,x_1]$ be a complete excursion, and list the nonzero increments in
$I$ as
\[
 x_0<a_1<a_2<\cdots<a_L\le x_1,
 \qquad \mu(a_j)\ne0.
\]
Put
\[
 r_I(t):=\#\{j:a_j\le t\},\qquad x_0<t\le x_1.
\]
If $\sigma\in\{\pm1\}$ is the orientation of the excursion, define
\[
 S_0=0,\qquad
 S_r:=\sigma\sum_{j=1}^r\mu(a_j)\qquad(1\le r\le L).
\]
Then $L=2n$ is even and
\[
 S_0=S_L=0,\qquad S_r>0\quad(1\le r<L),\qquad
 S_r-S_{r-1}\in\{-1,+1\}.
\]
Thus $(S_r)_{0\le r\le L}$ is a primitive Dyck path.  The
\emph{fixed-mask primitive Catalan law} keeps the physical support
$\{a_1,\dots,a_L\}$ fixed and assigns to those support positions a uniformly
chosen primitive Dyck sign sequence of length $L$.  We denote its probability
by $\mathbb P_{\mathrm{pCat},I}$.
\end{definition}

For the observed excursion one has, exactly,
\begin{equation}
\label{eq:support-time-height}
 |M(t)|=S_{r_I(t)},
 \qquad
 r_I(t)\le t-x_0<t.
\end{equation}
The first equality uses the orientation of the excursion and the fact that
$M$ is constant across the zero increments.

We first record the global Catalan estimate, since it is useful in its own
right and clarifies why a different estimate is needed at the square
horizon.

\begin{lemma}[Catalan maximum tail]
\label{lem:catalan-max-tail}
Let $L=2n$ and let $H_{\max}$ be the maximum of a uniformly chosen Dyck path
of length $2n$.  For $1\le h\le n$,
\begin{equation}
\label{eq:catalan-max-tail}
 \mathbb P_{\mathrm{Cat},L}\{H_{\max}\ge h\}
 \le (n+1)\frac{\binom{2n}{n+h}}{\binom{2n}{n}}
 \le (n+1)\exp\!\left(-\frac{h^2}{n+h}\right)
 \le \left(\frac L2+1\right)e^{-h^2/L}.
\end{equation}
For a uniformly chosen primitive Dyck path of length $L=2n$,
\begin{equation}
\label{eq:primitive-catalan-max-tail}
 \mathbb P_{\mathrm{pCat},L}\{H_{\max}\ge h\}
 \le \frac L2\exp\!\left(-\frac{(h-1)^2}{L-2}\right),
 \qquad 2\le h\le n.
\end{equation}
\end{lemma}

\begin{proof}
There are $\mathrm{Cat}_n=(n+1)^{-1}\binom{2n}{n}$ Dyck paths of length $2n$ and
$\binom{2n}{n}$ unrestricted bridges.  Reflecting a bridge after its first
visit to level $h$ gives a bijection between bridges that reach $h$ and
unrestricted walks ending at $2h$, hence their number is
$\binom{2n}{n+h}$.  Therefore
\[
 \mathbb P_{\mathrm{Cat},2n}\{H_{\max}\ge h\}
 \le (n+1)\frac{\binom{2n}{n+h}}{\binom{2n}{n}}.
\]
Moreover
\[
 \frac{\binom{2n}{n+h}}{\binom{2n}{n}}
 =\prod_{j=1}^h\frac{n-j+1}{n+j},
\]
and $\log(1-x)\le-x$ gives
\[
 \log\frac{\binom{2n}{n+h}}{\binom{2n}{n}}
 \le-\sum_{j=1}^h\frac{2j-1}{n+j}
 \le-\frac{h^2}{n+h}.
\]
Since $n+h\le2n=L$, \eqref{eq:catalan-max-tail} follows.  A primitive path
of length $2n$ has the form $+1$, followed by a Dyck path of length $2n-2$,
followed by $-1$; its maximum is one plus the maximum of the interior Dyck
path.  Applying the preceding bound with $n-1$ and $h-1$ proves
\eqref{eq:primitive-catalan-max-tail}.
\end{proof}

\begin{remark}[Why the maximum tail is not the square-horizon estimate]
\label{rem:max-tail-square-limit}
The exponent in \eqref{eq:primitive-catalan-max-tail} is of order
$H^2/L$.  For a general family of disjoint intervals retained by the
square horizon one may have $L\asymp x_0^2$ (for $M$ itself such excursions
are finitely many by Corollary~\ref{cor:KP-local-returns}), while the relevant witness lies much earlier, at physical
time $t\asymp x_0$.  Then a violation of size $t^{1/2+\eps}$ need not make
$H^2/L$ large.  Thus the global maximum tail, although correct, is not the
right Catalan quantity for combining with the fractional-square horizon.
The correct quantity is the pointwise height at the support time
$r_I(t)$.
\end{remark}

\begin{lemma}[Pointwise primitive-Catalan tail]
\label{lem:pointwise-catalan-tail}
Let $L=2n\ge4$ and let $(S_r)$ be a uniformly chosen primitive Dyck path of
length $L$.  For $1\le r\le L-1$ and $h\ge1$,
\begin{equation}
\label{eq:pointwise-catalan-tail-left}
 \mathbb P_{\mathrm{pCat},L}\{S_r\ge h\}
 \le 8n^{3/2}\exp\!\left(-\frac{h^2}{2r}\right)
 \le 3L^{3/2}\exp\!\left(-\frac{h^2}{2r}\right).
\end{equation}
By time reversal the same estimate holds with $r$ replaced by $L-r$;
hence
\begin{equation}
\label{eq:pointwise-catalan-tail-two-sided}
 \mathbb P_{\mathrm{pCat},L}\{S_r\ge h\}
 \le 3L^{3/2}
 \exp\!\left[-\frac{h^2}{2\min\{r,L-r\}}\right].
\end{equation}
If $h>\min\{r,L-r\}$ the probability is $0$, so the bound is only
relevant for $h\le\min\{r,L-r\}$.
\end{lemma}

\begin{proof}
Ignore both positivity and the endpoint constraint.  Among all $2^L$ sign
words, the proportion whose first $r$ signs have sum at least $h$ is at most
$\exp(-h^2/(2r))$ by the elementary Chernoff bound for a Rademacher sum.
Therefore the number of primitive Dyck paths satisfying $S_r\ge h$ is at
most
\[
 2^L\exp\!\left(-\frac{h^2}{2r}\right).
\]
The total number of primitive Dyck paths of length $2n$ is $\mathrm{Cat}_{n-1}$.  The
standard central-binomial lower bound
$\binom{2m}{m}\ge4^m/(2\sqrt m)$ gives
\[
 \mathrm{Cat}_{n-1}=\frac1n\binom{2n-2}{n-1}
 \ge\frac{4^{n-1}}{2n\sqrt{n-1}}.
\]
Since $2^L=4^n$, division gives the first inequality in
\eqref{eq:pointwise-catalan-tail-left}; the second follows from
$8n^{3/2}<3(2n)^{3/2}$.  Reversing the path and changing the signs of the
reversed increments preserves the primitive Dyck class and carries height
$S_r$ to height $S_{L-r}$, which proves
\eqref{eq:pointwise-catalan-tail-two-sided}.
\end{proof}

For an observed excursion $I$ define its pointwise Catalan upper-tail value
by
\begin{equation}
\label{eq:pointwise-catalan-pvalue}
 p_{\mathrm{Cat}}(I,t):=
 \mathbb P_{\mathrm{pCat},I}
 \left\{S_{r_I(t)}\ge |M(t)|\right\}.
\end{equation}
This is a tail probability under the fixed-mask Catalan reference law; it is
not a probability assigned to the arithmetic M\"obius path itself.

\begin{corollary}[Fractional-square Catalan rarity]
\label{cor:fractional-square-catalan-rarity}
Fix $c>0$ and $\eps>0$.  Let $I=(x_0,x_1]$ be a complete excursion with
$x_1\le cx_0^2$, and keep its zero mask fixed.  For all sufficiently large
$x_0$,
\begin{equation}
\label{eq:fractional-square-catalan-union}
 \mathbb P_{\mathrm{pCat},I}
 \left\{\exists\,t\in I:\ S_{r_I(t)}\ge t^{1/2+\eps}\right\}
 \le 3c^{5/2}x_0^5
 \exp\!\left(-\frac12x_0^{2\eps}\right).
\end{equation}
In particular the right-hand side is smaller than $x_0^{-A}$ for every
fixed $A>0$ once $x_0$ is large enough.
\end{corollary}

\begin{proof}
At a support position $t=a_r$ one has $r\le t-x_0<t$.  Hence
Lemma~\ref{lem:pointwise-catalan-tail} with $h=t^{1/2+\eps}$ gives
\[
 \mathbb P_{\mathrm{pCat},I}\{S_r\ge t^{1/2+\eps}\}
 \le3L^{3/2}\exp\!\left(-\frac12t^{2\eps}\right)
 \le3L^{3/2}\exp\!\left(-\frac12x_0^{2\eps}\right).
\]
There are at most $L$ support positions.  Since
$L\le x_1-x_0\le x_1\le cx_0^2$, the union bound gives
\eqref{eq:fractional-square-catalan-union}.  Between support positions the
path is constant, so no additional physical times have to be counted.
\end{proof}

The preceding statement is purely combinatorial: it measures how far a large
value lies in the tail of a reference law, not how the M\"obius constraints
select a path. Combined with the localization theorem it gives the following
quantitative form.

\begin{corollary}[Catalan surprise of a violation]
\label{cor:catalan-surprise}
\label{thm:catalan-selector-square}
Suppose $\Theta>1/2$, fix $0<\eta<\Theta-\frac12$, put
$\delta=\Theta-\eta-\frac12$, and fix $\lambda>\lambda_*$. Then there are
infinitely many excursions $I=(x_0,x_1]$ with $x_1\le\lambda x_0$ containing
a point $t$ with
\begin{equation}
\label{eq:KP-catalan-surprise}
-\log p_{\mathrm{Cat}}(I,t)\ \ge\ \frac18x_0^{2\delta}-O(\log x_0).
\end{equation}
The same holds with $x_1\le\lambda x_0$ replaced by $x_1\le cx_0^2$ for any
fixed $c>0$.
\end{corollary}

\begin{proof}
By Corollary~\ref{cor:KP-large-excursions} and Lemma~\ref{lem:linear-lacunary},
at least $(\gamma_1/\pi-1/\log\lambda-o(1))\log X$ excursions starting at
most $X$ have $x_1\le\lambda x_0$ and contain $t$ with
$h=|M(t)|\ge\frac12t^{1/2+\delta}$; for $cx_0^2$ use
Lemma~\ref{lem:partition} instead. With $r=r_I(t)<t$ and $L\le x_1-x_0$,
Lemma~\ref{lem:pointwise-catalan-tail} gives
$p_{\mathrm{Cat}}(I,t)\le3L^{3/2}\exp(-h^2/(2r))\le3L^{3/2}\exp(-t^{2\delta}/8)$,
and $t>x_0$, $L\le\lambda x_0$ (respectively $L\le cx_0^2$).
\end{proof}

\begin{remark}[Selector bounds are pointwise bounds]
\label{rem:catalan-selector-gap}
Up to polynomial factors, $p_{\mathrm{Cat}}(I,t)$ is a decreasing function of
$|M(t)|^2/r_I(t)$: by Lemma~\ref{lem:pointwise-catalan-tail},
\[
-\log p_{\mathrm{Cat}}(I,t)\ \ge\ \frac{|M(t)|^2}{2r_I(t)}-\log\bigl(3L^{3/2}\bigr).
\]
Hence a selector bound $-\log p_{\mathrm{Cat}}(I,t)\le\Phi(x_0)$ with
$\Phi(x)=x^{o(1)}$ on the excursions with $x_1\le\lambda x_0$ implies
$|M(t)|\le t^{1/2}x_0^{o(1)}$ there, and RH then follows from
Corollary~\ref{cor:descending}; it is not an independent criterion. The
selector bound is strictly stronger than Littlewood's bound on these
ranges: it controls $|M(t)|$ at the diffusive scale $r_I(t)^{1/2}$ of the
support time, which may be much smaller than $t^{1/2}$, and a polynomial
selector $p_{\mathrm{Cat}}\ge x_0^{-A}$ would give $|M(t)|\ll\sqrt{t\log t}$
on the retained ranges, beyond any bound known to follow from RH. Turning
$p_{\mathrm{Cat}}$ into an exact measure of deviation would require the
matching lower bound $-\log p_{\mathrm{Cat}}(I,t)\ll|M(t)|^2/r_I(t)+\log L$
for the primitive Catalan tail, which is not proved here.
\end{remark}

\begin{remark}[The correct quantifier]
\label{rem:quantifier}
A single value $|M(t)|\ge t^{1/2+\eps}$, or many values with bounded ratio
$|M(t)|/t^{1/2+\eps}$, is compatible with RH, since Littlewood's criterion
allows a constant $C_\eps$. The negation of the criterion is
\[
\exists\,\eps_0>0:\qquad
\sup_{t\ge1}\frac{|M(t)|}{t^{1/2+\eps_0}}=\infty,
\]
equivalently there are $t_j\to\infty$ and $A_j\to\infty$ with
$|M(t_j)|\ge A_jt_j^{1/2+\eps_0}$. Accordingly, the rarity bound
\eqref{eq:fractional-square-catalan-union} is a statement about the
reference law, and the arithmetic content lies in the sequence form
below. If $\Theta>1/2$ and $0<\delta<\Theta-\frac12$, the oscillation
theorem gives $|M(t_j)|\ge\frac12t_j^{1/2+\delta}$ along a sequence, and then
$|M(t_j)|/t_j^{1/2+\eps}\ge\frac12t_j^{\delta-\eps}\to\infty$ for every
$0<\eps<\delta$.
\end{remark}

\begin{corollary}[Catalan cost along a Littlewood-violating sequence]
\label{cor:catalan-sequence}
Let $t_j\to\infty$, $A_j\to\infty$ and $\eps_0>0$ with
$|M(t_j)|\ge A_jt_j^{1/2+\eps_0}$, and let $I_j=(x_{0,j},x_{1,j}]$ be the
excursion containing $t_j$. If $x_{1,j}\le x_{0,j}^{B}$ for a fixed $B$,
then
\[
-\log p_{\mathrm{Cat}}(I_j,t_j)\ \ge\ \frac{A_j^2}{2}\,t_j^{2\eps_0}
-O_B(\log x_{0,j})\ \ge\ \frac{A_j^2}{2}\,x_{0,j}^{2\eps_0}-O_B(\log x_{0,j}).
\]
\end{corollary}

\begin{proof}
With $r=r_{I_j}(t_j)<t_j$ and $h=|M(t_j)|$,
$h^2/(2r)\ge A_j^2t_j^{1+2\eps_0}/(2t_j)$; apply
Lemma~\ref{lem:pointwise-catalan-tail} with $L\le x_{1,j}\le x_{0,j}^B$.
\end{proof}

\subsubsection*{Selection energies: complete excursions and open prefixes}

The support-time clock of Definition~\ref{def:support-catalan} is the
squarefree-compressed clock
\[
\Sq(x)=\#\{n\le x:\mu(n)\ne0\},\qquad L^\#(I)=\Sq(x_1)-\Sq(x_0)\le x_1-x_0,
\]
and $r_I(t)=\Sq(t)-\Sq(x_0)$. Two tail quantities are natural: one for the
height of a complete excursion, and one for an open prefix, which does not
require knowledge of the future return.

\begin{definition}[Selection energies]
\label{def:selection-energies}
For a complete excursion $I$ with height $H(I)$ and compressed length
$L^\#=L^\#(I)$, put
\[
p^{\max}_{\mathrm{Cat}}(I):=\mathbb P_{\mathrm{pCat},L^\#}\{H_{\max}\ge H(I)\},
\qquad
\mathcal I_{\mathrm{Cat}}(I):=-\log p^{\max}_{\mathrm{Cat}}(I).
\]
For $t\ge1$ with $M(t)\ne0$ let $a(t)$ be the last zero of $M$ before $t$ and
$r=\Sq(t)-\Sq(a(t))$. After orientation, the compressed segment from $a(t)$ to
$t$ is a positive meander of length $r$: a $\pm1$ path $(S_k)_{k\le r}$ with
$S_0=0$ and $S_k>0$ for $1\le k\le r$. Let $\mathcal M_r$ be the set of such
paths and
\[
p_{\mathrm{open}}(t):=\mathbb P_{\mathcal M_r}\{S_r\ge|M(t)|\},
\]
the uniform measure on $\mathcal M_r$.
\end{definition}

\begin{lemma}[Height--entropy inequalities]
\label{lem:height-entropy}
\begin{enumerate}[label=\textup{(\roman*)}]
\item If $L^\#=2$, the compressed path is $0,1,0$ and $H(I)=1$. For every
complete excursion with $L^\#\ge4$,
\[
\mathcal I_{\mathrm{Cat}}(I)\ \ge\ \frac{(H(I)-1)^2}{L^\#-2}-\log\frac{L^\#}2 .
\]
\item $\#\mathcal M_r=\binom{r-1}{\lfloor(r-1)/2\rfloor}\ge2^r/(4\sqrt r)$,
and for $h\ge1$
\[
\mathbb P_{\mathcal M_r}\{S_r\ge h\}\le4\sqrt r\,e^{-h^2/(2r)},
\qquad
-\log p_{\mathrm{open}}(t)\ \ge\ \frac{M(t)^2}{2r}-\log\bigl(4\sqrt r\bigr).
\]
\end{enumerate}
\end{lemma}

\begin{proof}
(i) is \eqref{eq:primitive-catalan-max-tail}. (ii) The count of positive
meanders is classical (ballot theorem; the number ending at $s$ is
$\frac sr\binom{r}{(r+s)/2}$). From
$\binom{2m}{m}\ge4^m/(2\sqrt m)$ and
$\binom{2m+1}{m}=\frac12\binom{2m+2}{m+1}$ one gets
$\#\mathcal M_r\ge2^r/(4\sqrt r)$. Among all $2^r$ words the proportion with
$S_r\ge h$ is at most $e^{-h^2/(2r)}$ by the Chernoff bound; divide.
\end{proof}

\begin{proposition}[Selector criteria]
\label{prop:selector-criteria}
\begin{enumerate}[label=\textup{(\roman*)}]
\item \emph{Complete excursions.} Fix $\lambda>\lambda_*$. If there is $A$
such that $p^{\max}_{\mathrm{Cat}}(I)\ge x_0^{-A}$ for every sufficiently
large complete excursion with $x_1\le\lambda x_0$, then
$H(I)\ll_{A,\lambda}\sqrt{x_0\log x_0}$ on that family, and RH holds.
\item \emph{Open prefixes.} If there is $A$ such that
$p_{\mathrm{open}}(t)\ge t^{-A}$ for all large $t$ with $M(t)\ne0$, then
$|M(t)|\ll_A\sqrt{t\log t}$, and RH holds.
\end{enumerate}
\end{proposition}

\begin{proof}
(i) If $L^\#=2$ then $H(I)=1$. Otherwise $L^\#\ge4$ and, by
Lemma~\ref{lem:height-entropy}(i), $(H-1)^2\le(L^\#-2)(A\log x_0+\log L^\#)$, and $L^\#\le(\lambda-1)x_0$.
Littlewood's bound therefore holds on every excursion with
$x_1\le\lambda x_0$, and Corollary~\ref{cor:descending}(iii), in the form
that discards the $\lambda$-long excursions, gives RH.
(ii) By Lemma~\ref{lem:height-entropy}(ii),
$M(t)^2\le2r(A\log t+\log(4\sqrt r))$ with $r\le t$.
\end{proof}

\begin{remark}[Status of the two selectors]
\label{rem:selector-status}
Both hypotheses are stronger than RH: each yields $|M|\ll\sqrt{t\log t}$ on
its range, beyond any bound known to follow from RH. Neither is known to follow
from RH alone: a ballistic stretch of compressed length $r$ has $-\log p\asymp r$,
and RH does not exclude such stretches with $r$ a small power of $t$. The open selector controls every point without reference
to the future return; the complete selector uses only the horizon-localized
family, where $L^\#\le(\lambda-1)x_0$ is controlled. By contraposition and
Corollary~\ref{cor:catalan-sequence}, failure of RH forces the compressed
M\"obius path to realize, along an unbounded sequence, sectors of Catalan or
meander weight $\exp(-c\,t^{2\delta})$ with $\delta>0$, not merely
polynomially rare ones.
\end{remark}





\begin{remark}[The threshold $\lambda_*$]
\label{rem:lambda-star}
The function $\cos(\gamma_1\log x)$ changes sign $\gamma_1/\pi$ times per
unit of $\log x$, i.e.\ once in every multiplicative interval of ratio
$e^{\pi/\gamma_1}=\lambda_*$; the constant in \eqref{eq:KP-signchanges} is
this density. On the geometric side, a chain of $\lambda$-long intervals
realizes the bound of Lemma~\ref{lem:linear-lacunary}. The method of
Theorem~\ref{thm:horizon} compares these two counts, and $\lambda_*$ is the
exact point where they balance: for $\lambda\le\lambda_*$ the $\lambda$-long
family can, as far as counting alone is concerned, absorb every
Kaczorowski--Pintz excursion. Going below $\lambda_*$ requires information
on where inside an excursion the large values occur, not only how many
excursions carry them.
\end{remark}

\begin{proposition}[Logarithmic-length budget]
\label{prop:log-budget}
Let $\F$ be a family of excursions and let $u_1<\dots<u_N$ be the starts of
its members with $2\le u_j\le X$, with ends $v_j$. Then
\[
\sum_{j=1}^{N-1}\log\frac{v_j}{u_j}\le\log\frac X2 .
\]
If moreover $v_N\le u_N^A$ for some $A\ge1$, then
$\sum_{j=1}^{N}\log(v_j/u_j)\le A\log X$.
\end{proposition}

\begin{proof}
Disjointness gives $u_{j+1}\ge v_j$, so
$\log(v_j/u_j)\le\log u_{j+1}-\log u_j$; the sum telescopes to at most
$\log u_N-\log u_1\le\log(X/2)$. The last term is at most
$(A-1)\log u_N\le(A-1)\log X$.
\end{proof}

Lemmas~\ref{lem:power-lacunary}--\ref{lem:linear-lacunary} are the special
cases in which every term is bounded below. The budget form shows what the
horizon theorem actually measures.

\begin{corollary}[Mean logarithmic length of Kaczorowski--Pintz excursions]
\label{cor:KP-mean-log-length}
Suppose $\Theta>1/2$, fix $0<\eta<\Theta-\frac12$, and let $\mathcal K(X)$
be any set of distinct excursions starting at most $X$, each containing a
point $t$ with $|M(t)|\ge\frac12t^{\Theta-\eta}$, such that
$\#\mathcal K(X)\ge(\gamma_1/\pi-o(1))\log X$ (such sets exist by
Corollary~\ref{cor:KP-large-excursions}). Then
\[
\frac1{\#\mathcal K(X)}\sum_{\nu\in\mathcal K(X)}\log\frac{v_\nu}{u_\nu}
\ \le\ \frac{\pi}{\gamma_1}+o(1)=\log\lambda_*+o(1).
\]
\end{corollary}

\begin{proof}
Fix $A>1$. By Corollary~\ref{cor:KP-local-returns}, all but finitely many
excursions satisfy $v\le u^A$; the finitely many others contribute $O_A(1)$.
Proposition~\ref{prop:log-budget} gives a total of at most
$A\log X+O_A(1)$. Divide by $\#\mathcal K(X)$ and let $A\downarrow1$.
\end{proof}

\begin{remark}[What lowering $\lambda_*$ requires]
\label{rem:surplus}
Corollary~\ref{cor:descending}(iii) is the statement that the
excursions carrying Kaczorowski--Pintz witnesses cannot all have
$\log(v/u)>\log\lambda$ with $\lambda>\lambda_*$;
Corollary~\ref{cor:KP-mean-log-length} is its averaged form. Consequently,
any proof that such excursions have mean logarithmic length at least
$\pi/\gamma_1+\delta$ for some fixed $\delta>0$ would contradict
Corollary~\ref{cor:KP-mean-log-length} and prove RH, and a surplus
$\delta>\log(\lambda_*/\lambda)$ among the excursions whose witness lies
beyond $\lambda x_0$ would extend Corollary~\ref{cor:descending}(iii) to
that $\lambda<\lambda_*$. Kaczorowski--Pintz, as stated, gives neither: it
supplies at least one return in each band $[Y^{1-\delta},Y]$ but no upper
separation between consecutive witnesses. Nor do short-interval estimates
for $\mu$ supply it: $|M(x+h)-M(x)|=o(h)$ is compatible with $M$ staying
positive throughout $(x,x+h]$, whereas a return requires a decrement of
size at least $M(x)$ in a prescribed direction. Energy bounds do not help
either, since a flat excursion of height $1$ has $\beta_I,Q_I\asymp\ell_I$.
\end{remark}


\subsection{The linear window and the one-step recursion}
\label{subsec:linear-window}

Once the horizon is below $2x_0$, the square-range bilinear structure of
Proposition~\ref{prop:prefix-bilinear} is no longer needed: on $(u,2u]$ the
M\"obius function is a \emph{linear} expression in the prefix
$\mu\mathbf 1_{[1,u]}$.

\begin{proposition}[Linear window]
\label{prop:linear-window}
Let $u\ge1$. For every integer $u<n\le2u$,
\begin{equation}
\label{eq:linear-window-mu}
\mu(n)=-\sum_{\substack{d\mid n\\ d\le u}}\mu(d).
\end{equation}
Consequently, if $M(u)=0$ and $u<t\le2u$,
\begin{equation}
\label{eq:linear-window-M}
M(t)=-\sum_{d\le u}\mu(d)\left(\left\lfloor\frac td\right\rfloor
-\left\lfloor\frac ud\right\rfloor\right).
\end{equation}
\end{proposition}

\begin{proof}
Since $n\ge2$, $\sum_{d\mid n}\mu(d)=0$. Every proper divisor of $n$ is at
most $n/2\le u$, and every divisor $d\le u$ is proper because $n>u$. Summing
\eqref{eq:linear-window-mu} over $u<n\le t$, the number of multiples of $d$
in $(u,t]$ is $\lfloor t/d\rfloor-\lfloor u/d\rfloor$.
\end{proof}

\begin{corollary}[Localized prefix formula]
\label{cor:localized-prefix}
Let $M(u)=0$, let $1<\lambda<2$, and let $u<t\le\lambda u$.  Then
\begin{equation}
\label{eq:localized-prefix}
M(t)=-\sum_{d\le\lambda u/2}\mu(d)
\left(\left\lfloor\frac td\right\rfloor-
\left\lfloor\frac ud\right\rfloor\right).
\end{equation}
Indeed, if $\lambda u/2<d\le u$, then $d>u/2$ and the first multiple of
$d$ strictly larger than $u$ is $2d>\lambda u\ge t$.
\end{corollary}

\begin{corollary}[Prefix-dependency cutoff]
\label{cor:prefix-cutoff}
Put
\[
\theta_*:=\frac12e^{\pi/\gamma_1}=0.6244\ldots .
\]
For every fixed $\theta$ with $\theta_*<\theta<1$, RH is equivalent to
Littlewood's bound on every excursion starting at $u$ throughout
\[
u<t\le\min\{x_1,2\theta u\},
\]
where, throughout that window, $M(t)$ is given by the sum
\eqref{eq:localized-prefix} using only the prefix
$\mu(d)$ with $d\le\theta u$.  In particular, the choice
$\theta=5/8$ corresponds to $\lambda=5/4$ and gives the RH-equivalent
window
\[
1\le r=t-u\le\min\{\ell-1,u/4\}
\]
using only $d\le5u/8$.

The number $\theta_*$ is an infimum, not an attained endpoint, because the
horizon theorem requires $\lambda>\lambda_*$.
\end{corollary}

\begin{proof}
Take $\lambda=2\theta>\lambda_*$ in
Corollary~\ref{cor:descending}(iii), then use
Corollary~\ref{cor:localized-prefix}.
\end{proof}


\begin{corollary}[Algebraic dependency gap]
\label{cor:dependency-gap}
Let $\theta_*<\theta<1$.  If $u$ is a zero of $M$, then every increment
$\mu(n)$ with
\[
u<n\le2\theta u
\]
is determined by the prefix
$\mu(1),\ldots,\mu(\lfloor\theta u\rfloor)$.  Consequently the intermediate
block
\[
(\theta u,u]
\]
does not occur as a direct input in the divisor recursion determining the
localized future window.  As $\theta\downarrow\theta_*$, the omitted block
has relative length tending to
$1-\theta_*=0.3755\ldots$.
\end{corollary}

\begin{proof}
For $u<n\le2\theta u$, every proper divisor $d\mid n$ satisfies
$d\le n/2\le\theta u$.  Hence
\[
\mu(n)=-\sum_{\substack{d\mid n\\ d<n}}\mu(d)
\]
uses only $d\le\theta u$.
\end{proof}

\begin{remark}[Direct dependency is not independence]
Corollary~\ref{cor:dependency-gap} is an algebraic statement about which
values occur as inputs of the divisor recursion.  It does not assert
statistical or arithmetic independence between the omitted block
$(\theta u,u]$ and the earlier prefix.
\end{remark}

In convolution form, \eqref{eq:linear-window-mu} says that
$(\mu-\mu_u)*1=\mu\mathbf 1_{(u,\infty)}$ on $n\le2u$, the order-one analogue
of \eqref{eq:bilinear-general}, which for $u<n\le u^2$ reads
$\mu(n)=-(\mu_u*\mu_u*1)(n)$; on $(u,2u]$ both give $\mu(n)$, the
bilinear form being the quadratic continuation of the linear one. Equivalently, \eqref{eq:linear-window-M}
follows from $\sum_{d\le t}\mu(d)\lfloor t/d\rfloor=1$, since
$\lfloor t/d\rfloor=1$ for $u<d\le t\le2u$.

\begin{proposition}[One-step recursion]
\label{prop:one-step}
For real $x\ge1$,
\begin{equation}
\label{eq:hyperbola-M}
\sum_{k\le x}M\!\left(\frac xk\right)=1 .
\end{equation}
Hence, if $M(u)=0$ and $t>u$,
\begin{equation}
\label{eq:descent-identity}
M(t)=-\sum_{2\le k\le t}\left[M\!\left(\frac tk\right)-M\!\left(\frac uk\right)\right],
\end{equation}
with $M(y)=0$ for $y<1$. Every argument on the right is at most $t/2$; hence on a localized
window $u<t\le\lambda u$ all arguments are at most $\lambda u/2$.  For
$\lambda=5/4$ this is $5u/8$.
\end{proposition}

\begin{proof}
$\sum_{k\le x}M(x/k)=\sum_{k\le x}\sum_{d\le x/k}\mu(d)
=\sum_{m\le x}\sum_{d\mid m}\mu(d)=1$. Subtract the identity at $u$, where
$M(u)=0$, and note $M(u/k)=0$ for $k>u$.
\end{proof}

\begin{proposition}[Scale cascade in the localized excursion window]
\label{prop:scale-cascade}
Let $(u,v]$ be an excursion, let $\lambda_*<\lambda<2$, and let
$u<t\le\min\{v,\lambda u\}$.  For $k\ge2$ put
\[
\Delta_k(u,t):=M(t/k)-M(u/k),
\]
and for an integer $K\ge2$ put
\[
\mathcal R_K(u,t):=-\sum_{K<k\le t}\Delta_k(u,t).
\]
Then, exactly,
\begin{equation}
\label{eq:scale-cascade}
M(t)=-\sum_{k=2}^{K}\Delta_k(u,t)+\mathcal R_K(u,t),
\end{equation}
and $\mathcal R_K(u,t)$ uses only values $M(q)$ with
\[
q\le\frac{\lambda u}{K+1}.
\]
Equivalently, all dependence on values above $\lambda u/(K+1)$ is
concentrated in the $K-1$ increments
$\Delta_2,\ldots,\Delta_K$.  Moreover
\[
\Delta_k(u,t)=\sum_{u/k<n\le t/k}\mu(n),
\]
so the supporting interval has length $(t-u)/k$.
\end{proposition}

\begin{proof}
Equation~\eqref{eq:scale-cascade} is Proposition~\ref{prop:one-step} with
the terms $2\le k\le K$ separated.  If $k>K$, then
$t/k\le\lambda u/(K+1)$ and $u/k\le\lambda u/(K+1)$.  The last identity is
the definition of the summatory function on the interval $(u/k,t/k]$.
\end{proof}

\begin{corollary}[First channels]
\label{cor:first-channels}
For $K=2$,
\[
M(t)=-\bigl(M(t/2)-M(u/2)\bigr)+R_2(u,t),
\]
where $R_2$ uses only $q\le\lambda u/3$.  With the convenient localized
choice $\lambda=5/4$, the only channel above $5u/12$ is
\[
\Delta_2(u,u+r)=\sum_{u/2<n\le u/2+r/2}\mu(n),
\qquad 0<r\le u/4,
\]
whose supporting interval has length at most $u/8$.  In the limiting
sense $\lambda\downarrow\lambda_*$, the successive remainder cutoffs are
\[
\frac{\lambda_*}{2}u=0.6244\ldots u,\quad
\frac{\lambda_*}{3}u=0.4162\ldots u,\quad
\frac{\lambda_*}{4}u=0.3122\ldots u,\quad
\frac{\lambda_*}{5}u=0.2497\ldots u.
\]
\end{corollary}

\begin{lemma}[Channel persistence extracted from a high excursion]
\label{lem:channel-persistence}
Let $S_r=\sigma M(u+r)>0$ on $1\le r\le h$ with $S_0=0$ and
$|S_r-S_{r-1}|\le1$.  Fix $K\ge2$ and write, using
Proposition~\ref{prop:scale-cascade},
\[
S_r=\sum_{k=2}^{K}X_k(r)+Y_K(r),
\quad
X_k(r):=-\sigma\Delta_k(u,u+r),
\quad
Y_K(r):=\sigma \mathcal R_K(u,u+r).
\]
If $H=\max_{r\le h}S_r$ and $\tau$ is a point with $S_\tau=H$, then among
the last $\lfloor H/2\rfloor+1$ indices
\[
\tau-\lfloor H/2\rfloor,\ldots,\tau
\]
there is one fixed component among
$X_2,\ldots,X_K,Y_K$ which is at least $H/(2K)$ at no fewer than
$\lfloor H/(2K)\rfloor$ indices.
\end{lemma}

\begin{proof}
Since $S_0=0$ and the path has unit steps, $\tau\ge H$ and
$S_{\tau-j}\ge H-j\ge H/2$ for
$0\le j\le\lfloor H/2\rfloor$.  At every one of these indices the sum of
the $K$ displayed components is at least $H/2$, hence some component is at
least $H/(2K)$.  Pigeonholing the indices among the $K$ components gives the
claim.
\end{proof}

\begin{remark}[What the persistence lemma adds]
\label{rem:persistence-content}
Lemma~\ref{lem:channel-persistence} is not a bound for $M$.  It converts a
large value attained inside one excursion into a quantified persistence
requirement: either one of the finitely many scaled interval sums
$\Delta_2,\ldots,\Delta_K$ repeatedly carries a positive fraction of the
height, or the responsibility lies repeatedly in the deeper remainder,
which is supported entirely below $\lambda u/(K+1)$.  This avoids the
iteration defect of Remark~\ref{rem:iteration-gap}: no new point is claimed
to be the start-window point of another excursion.
\end{remark}
\begin{corollary}[Channel descent with explicit loss]
\label{cor:channel-descent}
Let $u$ be a zero of $M$, $t>u$, $K\ge2$, and $0<\alpha<1$. Then either
\begin{enumerate}[label=\textup{(\roman*)}]
\item there are $2\le k\le K$ and $q\in\{\lfloor u/k\rfloor,\lfloor t/k\rfloor\}$
with $|M(q)|\ge|M(t)|/(2K)$, hence
\[
\frac{|M(q)|}{q^\alpha}\ge\frac{k^\alpha}{2K}\cdot\frac{|M(t)|}{t^\alpha},
\qquad q\le\frac tk;
\]
or
\item $|\mathcal R_K(u,t)|\ge|M(t)|/K$, where $\mathcal R_K$ depends only on
$M(q)$ with $q\le t/(K+1)$.
\end{enumerate}
The statement holds at every $t>u$, not only in a localized window.
\end{corollary}

\begin{proof}
The proof of Proposition~\ref{prop:scale-cascade} applies verbatim to every
$t>u$, with $\mathcal R_K$ supported on $q\le t/(K+1)$. By \eqref{eq:scale-cascade}, one of the $K$ terms
$\Delta_2,\dots,\Delta_K,\mathcal R_K$ has modulus at least $|M(t)|/K$. If it is
$\Delta_k=M(\lfloor t/k\rfloor)-M(\lfloor u/k\rfloor)$, one of its two values
has modulus at least $|M(t)|/(2K)$, and $q\le t/k$.
\end{proof}

\begin{remark}[Channels do not iterate to a bound]
\label{rem:channel-loss}
Corollary~\ref{cor:channel-descent} has no iteration gap: it applies at
the new point $q$ relative to the last zero of $M$ before $q$. Its defect
is the constant loss $k^\alpha/(2K)<1$ per step. If alternative~(i) held at
every step, a chain of length $m\asymp\log t$ would give only
$|M(t)|\le(2K)^m\,O(1)=t^{O(\log K)}$, weaker than the trivial bound.
A closing argument must therefore show that alternative~(ii), or
alternative~(i) with loss factor at least $1$, occurs along the chain; both
are statements about cancellation, not about dependency.
\end{remark}


\begin{remark}[What the recursion determines]
\label{rem:recursion-content}
Iterating \eqref{eq:hyperbola-M} determines all of $M$ from $M(1)=1$, so
determinacy alone carries no information; cf.\
Remark~\ref{rem:information}.  The content of
\eqref{eq:descent-identity}, after Corollary~\ref{cor:prefix-cutoff}, is
more precise: at the start $u$ of an excursion, every value in an
RH-equivalent window $u<t\le\lambda u$ is a signed sum of increments of
$M$ over scaled intervals $(u/k,t/k]$ lying below $\lambda u/2$.  This is a
descent of dependency only; no descent of normalized size follows from the
identity.
\end{remark}

\begin{corollary}[RH in the localized linear window]
\label{cor:linear-window-RH}
RH holds if and only if, for every $\eps>0$,
\[
\left|\sum_{d\le5u/8}\mu(d)\left(\left\lfloor\frac td\right\rfloor
-\left\lfloor\frac ud\right\rfloor\right)\right|
=\left|\sum_{2\le k\le t}\left[M\left(\frac tk\right)-M\left(\frac uk\right)\right]\right|
\ll_\eps t^{1/2+\eps}
\]
for every excursion starting at $u$ and every
\[
u<t\le\min\left\{x_1,\frac54u\right\}.
\]
Equivalently, in intra-excursion coordinates $t=u+r$ it is enough to
control
\[
1\le r\le\min\{\ell-1,u/4\}.
\]
\end{corollary}

\begin{proof}
Use Corollary~\ref{cor:prefix-cutoff} with $\theta=5/8$, together with
Proposition~\ref{prop:one-step}.
\end{proof}

\begin{proposition}[Kernel form of the recursion]
\label{prop:kernel}
Let $M(u)=0$ and let $t$ be an integer with $u<t<2u$. For $1\le q\le t/2$ put
\[
c_q(t,u)=\Bigl(\Bigl\lfloor\frac tq\Bigr\rfloor-\Bigl\lfloor\frac t{q+1}\Bigr\rfloor\Bigr)
-\Bigl(\Bigl\lfloor\frac uq\Bigr\rfloor-\Bigl\lfloor\frac u{q+1}\Bigr\rfloor\Bigr).
\]
Then:
\begin{enumerate}[label=\textup{(\roman*)}]
\item $M(t)=-\sum_{1\le q\le t/2}c_q(t,u)M(q)$;
\item $c_q(t,u)=\dfrac{t-u}{q(q+1)}+e_q$ with $|e_q|<2$; for $q>\sqrt t$
one has $c_q\in\{-1,0,1\}$, and $c_q\ne0$ only if $q=\lfloor t/k\rfloor$ or
$q=\lfloor u/k\rfloor$ for some $k<\sqrt t$;
\item $c_1(t,u)\ge\frac12(t-u-1)$;
\item $M(t)=-(t-u)\,\Upsilon(\sqrt t)+\mathrm{Err}(t,u)$, where $\Upsilon$ is the
coefficient \eqref{eq:Theta_def} and
\[
|\mathrm{Err}(t,u)|\le2\sum_{q\le\sqrt t}|M(q)|
+\sum_{k<\sqrt t}\Bigl(\Bigl|M\Bigl(\frac tk\Bigr)\Bigr|+\Bigl|M\Bigl(\frac uk\Bigr)\Bigr|\Bigr).
\]
\end{enumerate}
\end{proposition}

\begin{proof}
(i) In \eqref{eq:descent-identity} group the indices $k\ge2$ by the value
$q=\lfloor t/k\rfloor$, respectively $q=\lfloor u/k\rfloor$. The number of
$k\ge1$ with $\lfloor t/k\rfloor=q$ is
$\lfloor t/q\rfloor-\lfloor t/(q+1)\rfloor$, and for $q\le t/2$ none of them
is $k=1$. The same holds for $u$, because $k=1$ gives $q=u>t/2$. (ii) Each
bracket equals $x/(q(q+1))$ minus a difference of two fractional parts, so
it differs from $x/(q(q+1))$ by less than $1$. For $q>\sqrt t\ge\sqrt u$ both
$t/q-t/(q+1)$ and $u/q-u/(q+1)$ are $<1$, so each bracket is $0$ or $1$, and
it is $1$ only if $q=\lfloor x/k\rfloor$ with $k\le x/q<\sqrt t$.
(iii) $c_1=(t-u)-(\lfloor t/2\rfloor-\lfloor u/2\rfloor)$ and
$\lfloor t/2\rfloor-\lfloor u/2\rfloor\le\lceil(t-u)/2\rceil$.
(iv) Insert (ii) into (i) for $q\le\sqrt t$; the smooth part is
$(t-u)\sum_{q\le\sqrt t}M(q)/(q(q+1))=(t-u)\Upsilon(\sqrt t)$. For
$q>\sqrt t$ use the last assertion of (ii).
\end{proof}

\begin{corollary}[Loss of the pointwise induction]
\label{cor:induction-loss}
Let $0<\alpha<1$ and suppose $|M(q)|\le Cq^\alpha$ for all $q\le t/2$. Then,
for $u<t<2u$ with $M(u)=0$,
\begin{equation}
\label{eq:induction-loss}
\bigl|M(t)+(t-u)\Upsilon(\sqrt t)\bigr|\ll_\alpha C\,t^{(1+\alpha)/2},
\end{equation}
while the termwise bound of the full kernel gives only
$|M(t)|\ll_\alpha C\,t$. For $\alpha=\frac12+\eps$ the pointwise induction
therefore loses $t^{1/2-\eps}$, and $t^{1/4-\eps/2}$ after the smooth
component $-(t-u)\Upsilon(\sqrt t)$ is extracted.
\end{corollary}

\begin{proof}
By Proposition~\ref{prop:kernel}(iv),
$|\mathrm{Err}|\le2C\sum_{q\le\sqrt t}q^\alpha+2Ct^\alpha\sum_{k<\sqrt t}k^{-\alpha}
\ll_\alpha Ct^{(1+\alpha)/2}$. The full kernel satisfies
$\sum_q|c_q|q^\alpha\ge c_1\ge\frac12(t-u-1)$ by (iii), and
$\sum_q|c_q|q^\alpha\ll_\alpha t$ by (ii).
\end{proof}

The smooth component is not small for termwise reasons: by
\eqref{eq:C_tail},
\[
\Upsilon(\sqrt t)=-\sum_{n>\sqrt t}\frac{M(n)}{n(n+1)},
\]
so its size is governed by the global identity
$\sum_{n\ge1}M(n)/(n(n+1))=0$, i.e.\ by the pole of $\zeta$ at $s=1$
(Remark~\ref{rem:why_log2}), and involves values of $M$ beyond $t/2$. Even
under RH the tail gives only $(t-u)\Upsilon(\sqrt t)\ll t^{3/4+\eps/2}$, the
same order as the bound for $\mathrm{Err}(t,u)$; the RH scale for $M(t)$ is a statement
about their sum.

\begin{remark}[Localized intra-excursion system]
\label{rem:sign-cone}
Let $I=(u,x_1]$ be an excursion of length $\ell=x_1-u$ and orientation
$\sigma$, and put
\[
S_r:=\sigma M(u+r)>0\qquad(1\le r<\ell),\qquad S_0=S_\ell=0.
\]
For a fixed $\lambda\in(\lambda_*,2)$ set
\[
h_\lambda(u):=\min\{\ell-1,(\lambda-1)u\}.
\]
Then the entire RH-relevant initial segment satisfies simultaneously
\[
S_r=-\sigma\sum_{q\le\lambda u/2}c_q(u+r,u)M(q)>0,
\qquad1\le r\le h_\lambda(u),
\]
and $|S_r-S_{r-1}|\le1$.  For $\lambda=5/4$ this is the segment
$r\le\min\{\ell-1,u/4\}$ and all dependencies lie below $5u/8$.
\end{remark}

\begin{definition}[Relaxed divisor-sum window class]
\label{def:relaxed-divisor-cone}
Fix $u\ge2$ and $1\le h<u$.  For
$\delta=(\delta_d)_{d\le u}\in\{-1,0,1\}^u$ with $\delta_1=1$ and
$\sum_{d\le u}\delta_d=0$, put
\[
D_\delta(n):=\sum_{\substack{d\mid n\\ d\le u}}\delta_d,
\qquad u<n\le u+h.
\]
We call $\delta$ window-admissible with orientation $\sigma\in\{-1,1\}$ if
$D_\delta(n)\in\{-1,0,1\}$ throughout the window and
\[
\sigma\sum_{u<n\le u+r}D_\delta(n)<0
\qquad(1\le r\le h).
\]
The induced oriented path is
$S_r=-\sigma\sum_{u<n\le u+r}D_\delta(n)>0$.
For the genuine M\"obius prefix, $\delta_d=\mu(d)$, these are exactly the
increment and excursion conditions coming from $\mu*1=\delta_1$.
\end{definition}

\begin{proposition}[Ballistic counterexample to the relaxed window class]
\label{prop:relaxed-ballistic}
The relaxed class of Definition~\ref{def:relaxed-divisor-cone} admits
ballistic paths.  Namely, take
\[
\delta_1=1,\qquad\delta_u=-1,\qquad\delta_d=0\quad(1<d<u).
\]
Then for every $u<n<2u$ one has $D_\delta(n)=1$, and with $\sigma=-1$,
\[
S_r=r\qquad(1\le r<u).
\]
If $m_q:=\sum_{d\le q}\delta_d$, then $m_q=1$ for every $q<u$.  Hence for
any $0<\alpha<1$, any fixed $1<\lambda<2$, and
$h=\lfloor(\lambda-1)u\rfloor$,
\[
\sup_{q\le\lambda u/2}\frac{|m_q|}{q^\alpha}=1,
\qquad
\frac{S_h}{(u+h)^\alpha}\asymp_\lambda u^{1-\alpha}\longrightarrow\infty.
\]
Thus a maximum principle comparing the normalized localized-window height
to the normalized history is false on this relaxed class.
\end{proposition}

\begin{proof}
For $u<n<2u$, the integer $u$ does not divide $n$.  Hence among the two
nonzero coordinates of $\delta$ only $d=1$ contributes to $D_\delta(n)$,
so $D_\delta(n)=1$.  The remaining assertions are immediate.
\end{proof}

\begin{remark}[What the relaxation loses]
\label{rem:relaxed-loses}
The counterexample violates the internal divisor identities of the M\"obius
prefix.  Already for every prime $p<u$ the genuine identity requires
$\delta_1+\delta_p=0$, whereas the ballistic vector above has
$\delta_p=0$.  Consequently any successful localized argument of this type
must use some of the relations
\[
\sum_{d\mid n}\delta_d=0\qquad(2\le n\le u),
\]
not merely the future-window constraints.
\end{remark}

\begin{remark}[The prefix correlation identity]
\label{rem:prefix-correlation}
The complete-excursion bridge identity must not be applied to a truncated
localized window.  For $1\le h<\ell$, put
$x_j=\sigma\mu(u+j)$,
$S_r=\sum_{j\le r}x_j$, and
$C_a^{[h]}=\sum_{j=1}^{h-a}x_jx_{j+a}$.  Lemma~\ref{lem:triangular-corr}
applied to the prefix of length $h$ gives exactly
\[
\sum_{r=1}^{h-1}S_r^2+\sum_{a=1}^{h-1}aC_a^{[h]}
=S_h\sum_{r=1}^{h-1}S_r.
\]
The right-hand side vanishes only for a complete bridge.  Thus
$Q_\nu=-\sum_a aC_a^{(\nu)}$ cannot be used as an identity internal to a
truncated horizon unless the excursion actually returns there.  In
particular, passing from a large localized value to the energy of the whole
excursion mixes the retained window with the part deliberately discarded by
the horizon theorem.
\end{remark}

\begin{remark}[The one-step iteration gap]
\label{rem:iteration-gap}
Even a one-step implication sending a large value in the localized window
to a large value at some $q\le\lambda u/2$ would not automatically iterate.
The point $q$ need not lie in the localized window of the excursion that
contains $q$; it may lie beyond $\lambda u'$ relative to that excursion's
start $u'$.  Sparsity of the family of $\lambda$-long excursions does not
repair a single descent chain.  Iteration therefore requires either a
principle valid for all $t$, or a mechanism that places the new point inside
an admissible window at every step.
\end{remark}

\begin{proposition}[Termwise absolute-value contraction gives no nontrivial gain]
\label{prop:cone-obstruction}
Let $(u,x_1]$ be an excursion, $1\le h<\min\{x_1-u,u\}$, and let
$w_1,\dots,w_h\ge0$, not all zero. Put
$K_q=\sum_{r\le h}w_rc_q(u+r,u)$, so that
$\sum_rw_rM(u+r)=-\sum_qK_qM(q)$. If, for some $0<\alpha<1$ and
$C_\alpha>0$,
\[
\sum_q|K_q|q^\alpha\le C_\alpha\sum_{r\le h}w_r(u+r)^\alpha,
\]
then
\[
\sum_{r\le h}w_r|M(u+r)|
\le(1+2C_\alpha)\sum_{r\le h}w_r(u+r)^\alpha.
\]
\end{proposition}

\begin{proof}
Proposition~\ref{prop:kernel}(iii) gives
$c_1(u+r,u)\ge(r-1)/2\ge0$, hence
$\frac12\sum_rw_r(r-1)\le K_1\le
C_\alpha\sum_rw_r(u+r)^\alpha$. Since $M(u)=0$ and $M$ has unit steps,
$|M(u+r)|\le r$, and the stated inequality follows.
\end{proof}

\begin{remark}[No strict coercivity from the first coefficient]
\label{rem:c1-parity}
Writing $t=u+r$ gives the exact formula
\[
c_1(u+r,u)=
\begin{cases}
\lceil r/2\rceil,&u\ \text{even},\\
\lfloor r/2\rfloor,&u\ \text{odd}.
\end{cases}
\]
In particular $c_1(u+1,u)=0$ when $u$ is odd.  The preceding proposition
therefore rules out only this particular termwise absolute-value mechanism;
it does not rule out a global signed contraction or cancellation among the
coefficients $K_q$.
\end{remark}

\begin{definition}[Multiplicative layers of the prefix]
\label{def:multiplicative-layers}
Let $\Omega(n)$ denote the number of prime factors of $n$, counted with
multiplicity.  For $R\ge2$ and $k\ge1$, a vector
$(\delta_d)_{d\le R}$ satisfies the first $k$ M\"obius layers if
$\delta_1=1$ and
\[
\sum_{d\mid n}\delta_d=0
\qquad(2\le n\le R,\ \Omega(n)\le k).
\]
Thus layer $1$ imposes the identities at primes, layer $2$ adds prime
squares and products of two primes, and so on.
\end{definition}

\begin{proposition}[One-layer propagation into the localized future]
\label{prop:layer-propagation}
Suppose $(\delta_d)_{d\le R}$ satisfies the first $k$ M\"obius layers. Then
\[
\delta_d=\mu(d)\qquad(d\le R,\ \Omega(d)\le k).
\]
Moreover, for every integer $n$ with
\[
R<n\le2R,\qquad\Omega(n)\le k+1,
\]
one has
\[
-\sum_{\substack{d\mid n\\d\le R}}\delta_d=\mu(n).
\]
Thus the $k$-th multiplicative layer of the prefix determines exactly the
$(k+1)$-st multiplicative layer of the future increments in a window of
length at most $R$.
\end{proposition}

\begin{proof}
The first assertion follows recursively from
$\delta_n=-\sum_{d\mid n,\ d<n}\delta_d$: every proper divisor of an
integer with $\Omega(n)\le k$ has smaller $\Omega$.  For the second, every
proper divisor of $n\le2R$ is at most $n/2\le R$ and has
$\Omega(d)\le k$.  Hence the displayed sum equals
$-\sum_{d\mid n,\ d<n}\mu(d)=\mu(n)$.
\end{proof}

\begin{question}[Multiplicative-layer depth]
\label{q:layer-depth}
Fix $\theta\in(\theta_*,1)$ and the localized window $u<n\le2\theta u$.
Let $k=k(u)$ be a depth up to which the internal M\"obius divisor identities
are imposed below $R=\theta u$. What growth of $k(u)$, if any, suffices to
combine those identities with $M(u)=0$ and the unit-step and excursion-sign
constraints so as to force
$\max_{u<t\le2\theta u}|M(t)|\ll_\eps u^{1/2+\eps}$? A bounded depth would
give a compressed arithmetic mechanism; imposing all layers reconstructs the
M\"obius prefix. Proposition~\ref{prop:layer-propagation} makes the depth a
quantifiable intermediate invariant; the results below treat only fixed
$k$.
\end{question}
\begin{proposition}[Layer sums in the localized window]
\label{prop:layer-sums}
Fix $1<\lambda<2$ and $j\ge1$. As $u\to\infty$,
\begin{equation}
\label{eq:layer-sum}
\sum_{\substack{u<n\le\lambda u\\ \Omega(n)=j}}\mu(n)
=(-1)^j\,(1+o(1))\,\frac{(\lambda-1)\,u\,(\log\log u)^{j-1}}{(j-1)!\,\log u}.
\end{equation}
\end{proposition}

\begin{proof}
If $\Omega(n)=j$ and $\mu(n)\ne0$, then $n$ is squarefree with exactly $j$
prime factors and $\mu(n)=(-1)^j$. By Landau's theorem (see
\cite[Ch.~7]{MontgomeryVaughan2007}), the number $\pi_j^\flat(x)$ of such
$n\le x$ satisfies $\pi_j^\flat(x)\sim f_j(x):=x(\log\log x)^{j-1}/((j-1)!\log x)$.
Since $f_j(x)/x$ is slowly varying, $f_j(\lambda u)-f_j(u)=(\lambda-1+o(1))f_j(u)$.
\end{proof}

\begin{corollary}[Bounded layer depth]
\label{cor:bounded-depth}
Fix $k\ge1$, $\theta\in(\theta_*,1)$, $R=\theta u$, and the window
$u<n\le2\theta u$.
\begin{enumerate}[label=\textup{(\roman*)}]
\item The increments fixed by the first $k$ layers
(Proposition~\ref{prop:layer-propagation}), namely those with
$\Omega(n)\le k+1$, form a set of $o(u)$ integers of the window.
\item Their summed contribution over the window is
\[
\sum_{\substack{u<n\le2\theta u\\ \Omega(n)\le k+1}}\mu(n)
=(-1)^{k+1}(1+o(1))\frac{(2\theta-1)\,u\,(\log\log u)^{k}}{k!\,\log u},
\]
so it exceeds $u^{1-\eps}$ in modulus and has a fixed sign.
\item Let $\delta_d=\mu(d)$ for $\Omega(d)\le k$ and $\delta_d=0$ otherwise.
For squarefree $n$ in the window with $\omega(n)=w\ge k+2$,
\[
D_\delta(n)=\sum_{\substack{d\mid n\\ d\le R}}\delta_d
=\sum_{j=0}^{k}(-1)^j\binom wj=(-1)^k\binom{w-1}{k},
\]
which has modulus at least $k+1\ge2$. Thus the zero extension of a
$k$-layer prefix violates the unit-step constraint at every such $n$, i.e.\
at all but $o(u)$ squarefree integers of the window.
\end{enumerate}
\end{corollary}

\begin{proof}
(i) and (ii) follow from Proposition~\ref{prop:layer-sums} summed over
$j\le k+1$, the term $j=k+1$ dominating. For (iii), all divisors of $n$
other than $n$ are $\le n/2\le R$, and $n$ itself has $\Omega(n)=w>k$, so
$D_\delta(n)$ is the truncated alternating binomial sum, evaluated by
induction on $k$ via $\binom{w-1}{k}=\binom wk-\binom{w-1}{k-1}$. The
proportion of squarefree $n\le x$ with $\omega(n)\le k+1$ tends to $0$ by
Proposition~\ref{prop:layer-sums}.
\end{proof}

\begin{remark}[What fixed layer depth proves]
\label{rem:depth-content}
For every fixed $k$, Corollary~\ref{cor:bounded-depth} shows that the first
$k$ layers determine only a density-zero subset of the localized window,
while the signed sum of the determined increments has modulus
$(1+o(1))(2\theta-1)u(\log\log u)^k/(k!\log u)$, larger than $u^{1-\eps}$
for all large $u$. Hence, if the full sum on the window is
$O(u^{1/2+\eps})$, substantial cancellation must occur between the
increments determined by any fixed set of initial layers and the remaining
ones. The results proved here do not determine how the required depth grows
with $u$; in particular they do not show that $\asymp\log\log u$ layers
(the typical size of $\Omega$) are necessary or sufficient, which would
require estimates uniform in $k=k(u)$. Part~(iii) shows that the unit-step
constraint is far from vacuous on a $k$-layer relaxation: it forbids the
naive completion at almost every integer of the window, so the free
coordinates $\delta_d$ with $\Omega(d)>k$ are strongly constrained.
\end{remark}


\subsection{Conditional extreme-value consequence of excursion sparsity}
\label{subsec:extreme}

The counts of Section~\ref{subsec:geometric} bound the number of excursions
in a family, not their heights. They acquire a second function when combined
with a tail hypothesis on a normalized excursion statistic: they fix the
scale of the family maximum, with no independence assumption.

Let $(\Omega,\mathcal G,\Pr)$ be a probability space carrying a random
sequence of disjoint integer intervals $I_k=(x_{0,k},x_{1,k}]$ with
$x_{1,k}\le x_{0,k+1}$, and real random variables $Z_k$ (for instance a model
of the M\"obius walk with $Z_k=H_k/\sqrt{\ell_k}$). A family $\F$ is given by
a rule selecting indices from $(x_{0,k},x_{1,k})_k$, such as $\F_{\mathfrak h}$; let
$k_\F(1)<k_\F(2)<\cdots$ enumerate its members, and put
\[
Z_{\max,\F}(X):=\max\{Z_k:k\in\F,\ x_{0,k}\le X\}\qquad(\max\emptyset=-\infty).
\]
The lemmas of Section~\ref{subsec:geometric} hold for every $\omega$.

\begin{hypothesis}[$\mathrm{SG}_\varpi$]
\label{hyp:SG}
There are $C,\kappa_0,\varpi>0$ such that for every $j\ge1$ and $z\ge0$
\[
\Pr\bigl(k_\F(j)\text{ exists and }Z_{k_\F(j)}>z\bigr)\le Ce^{-\kappa_0z^\varpi}.
\]
For $\varpi=2$ this is a uniform sub-Gaussian tail, denoted (SG).
\end{hypothesis}

\begin{theorem}[Extreme scale from a count]
\label{thm:ev-union}
Assume $(\mathrm{SG}_\varpi)$ for $\F$ and suppose that $E_\F(X)\le \Lambda(X)$ for
every $\omega$, with $\Lambda(X)\ge1$ deterministic. Then
\begin{equation}
\label{eq:ev-union}
\Pr\bigl(Z_{\max,\F}(X)>z\bigr)\le \Lambda(X)\,Ce^{-\kappa_0z^\varpi}.
\end{equation}
Consequently, for every $\delta>0$,
\begin{equation}
\label{eq:ev-scale}
\Pr\Bigl(Z_{\max,\F}(X)>\Bigl(\frac{1+\delta}{\kappa_0}\log \Lambda(X)\Bigr)^{1/\varpi}\Bigr)
\le C\,\Lambda(X)^{-\delta},
\end{equation}
and $Z_{\max,\F}(X)=O_{\Pr}\bigl((\log \Lambda(X))^{1/\varpi}\bigr)$ as $\Lambda(X)\to\infty$.
\end{theorem}

\begin{proof}
The members of $\F$ starting at most $X$ are the first $E_\F(X)\le \Lambda(X)$
members of the enumeration, so
$\{Z_{\max,\F}(X)>z\}\subset\bigcup_{j\le \Lambda(X)}\{k_\F(j)\text{ exists},\
Z_{k_\F(j)}>z\}$, and \eqref{eq:ev-union} is the union bound. Substituting
$\kappa_0z^\varpi=(1+\delta)\log \Lambda(X)$ gives \eqref{eq:ev-scale}.
\end{proof}

\begin{corollary}[Triple-logarithmic scale on the square family]
\label{cor:ev-square}
Assume (SG) for the square-exceptional family $\F_{x^2}$ of the model. Then,
for $X\ge e^e$ and every $\delta>0$,
\begin{equation}
\label{eq:ev-square}
\Pr\bigl(Z_{\max,\square}(X)>z\bigr)\le4C(\log\log X)e^{-\kappa_0z^2},
\end{equation}
\begin{equation}
\label{eq:ev-square-b}
\Pr\Bigl(Z_{\max,\square}(X)>\sqrt{\tfrac{1+\delta}{\kappa_0}\log\log\log X}\Bigr)
\le4C(\log\log X)^{-\delta}.
\end{equation}
Thus $Z_{\max,\square}(X)=O_{\Pr}(\sqrt{\log\log\log X})$. The same holds for
the family $\F_{cx^2}$ of Lemma~\ref{lem:partition} with a constant depending
on $c$.
\end{corollary}

\begin{proof}
Use \eqref{eq:square-exception-count} in \eqref{eq:ev-union} and note that
$e^{-\kappa_0z^2}=(\log\log X)^{-(1+\delta)}$ at the stated $z$.
\end{proof}

The union bound gives convergence in probability; a Borel--Cantelli argument
along a sequence adapted to $L$ gives an almost sure bound with the same
constant.

\begin{theorem}[Almost sure extreme scale]
\label{thm:ev-as}
Assume $(\mathrm{SG}_\varpi)$ for $\F$ and $E_\F(X)\le \Lambda(X)$ for every $\omega$,
where $L$ is continuous, nondecreasing and unbounded. Then almost surely
\[
\limsup_{X\to\infty}\frac{Z_{\max,\F}(X)}{(\log \Lambda(X))^{1/\varpi}}\le \kappa_0^{-1/\varpi}.
\]
\end{theorem}

\begin{proof}
Fix $\delta>0$. For $j\ge j_0$ let $X_j$ be the least $X$ with $\Lambda(X)=2^j$.
Put $z_j=\bigl((1+\delta)(j+1)\log2/\kappa_0\bigr)^{1/\varpi}$ and
$A_j=\{Z_{\max,\F}(X_{j+1})>z_j\}$. By \eqref{eq:ev-union},
$\Pr(A_j)\le2^{j+1}Ce^{-(1+\delta)(j+1)\log2}=C2^{-(j+1)\delta}$, which is
summable, so almost surely $A_j$ occurs only finitely often. For
$X_j\le X\le X_{j+1}$, monotonicity of $Z_{\max,\F}$ in $X$ gives
$Z_{\max,\F}(X)\le z_j$ outside $A_j$, while $\log \Lambda(X)\ge j\log2$. Hence
the $\limsup$ is at most
$\lim_j\bigl((1+\delta)(j+1)/(\kappa_0j)\bigr)^{1/\varpi}=((1+\delta)/\kappa_0)^{1/\varpi}$. Let
$\delta\downarrow0$ along a countable sequence.
\end{proof}

With $\Lambda_\square(X)=2+\log_2(\log X/\log2)$, which is continuous and
increasing, Theorem~\ref{thm:ev-as} gives under (SG)
\[
\limsup_{X\to\infty}\frac{Z_{\max,\square}(X)}{\sqrt{\log\log\log X}}\le \kappa_0^{-1/2}
\quad\text{almost surely.}
\]
Applied to the three counts available for any sequence of disjoint
intervals (all excursions: at most $X+1$; Lemma~\ref{lem:linear-lacunary};
Lemma~\ref{lem:power-lacunary}), Theorems~\ref{thm:ev-union}
and~\ref{thm:ev-as} give three levels:
\[
\begin{array}{c|c|c}
\text{family} & \text{members starting}\le X & \text{extreme scale under (SG)}\\
\hline
\text{all excursions} & \le X+1 & \sqrt{\log X}\\
x_1>\lambda x_0 & O(\log X) & \sqrt{\log\log X}\\
x_1>cx_0^2\ \ (\text{or }x_0^A) & O(\log\log X) & \sqrt{\log\log\log X}
\end{array}
\]
Each level of lacunarity removes one logarithm from the extreme scale.
Under $(\mathrm{SG}_\varpi)$ the last column reads $(\log X)^{1/\varpi}$,
$(\log\log X)^{1/\varpi}$, $(\log\log\log X)^{1/\varpi}$.

\begin{remark}[Scope]
\label{rem:ev-scope}
For $M$ itself the square-exceptional family is finite by
Corollary~\ref{cor:KP-local-returns}, so
Theorems~\ref{thm:ev-union}--\ref{thm:ev-as} concern stochastic models of
the M\"obius walk, in which the pathwise counts of
Section~\ref{subsec:geometric} persist while no analogue of the
Kaczorowski--Pintz localization is available. Hypothesis (SG) is an
assumption on the model; the theorems state exactly what the counts
contribute to its extremes.
\end{remark}

\subsection{The square range as an exact bilinear form of the prefix}

\begin{proposition}[Bilinear square-range identity]
\label{prop:prefix-bilinear}
Let $u\ge1$ and $1\le t\le u^2$. Then
\begin{equation}
\label{eq:bilinear-general}
M(t)=2M(\min(t,u))-
\sum_{d,j\le u}\mu(d)\mu(j)\left\lfloor\frac{t}{dj}\right\rfloor.
\end{equation}
In particular, if $M(u)=0$ and $u<t\le u^2$,
\begin{equation}
\label{eq:bilinear-excursion}
M(t)=-\mathcal B_u(t),\qquad
\mathcal B_u(t):=\sum_{d,j\le u}\mu(d)\mu(j)
\left\lfloor\frac{t}{dj}\right\rfloor,
\qquad \mathcal B_u(u)=0.
\end{equation}
\end{proposition}

\begin{proof}
Let $\mu_u=\mu\mathbf1_{[1,u]}$. Each factor of $(\mu-\mu_u)*(\mu-\mu_u)$ is
supported on integers $\ge u+1$, so the convolution, and its convolution
with $1$, vanish on $n<(u+1)^2$, in particular on $n\le u^2$. Expanding and using $\mu*1=\delta$ gives
$\mu(n)=2\mu_u(n)-(\mu_u*\mu_u*1)(n)$ for $n\le u^2$; sum over $n\le t$.
\end{proof}

\begin{remark}[Higher-order determinacy]
\label{rem:prefix-k}
For every $k\ge1$ and $n\le u^k$,
\begin{equation}
\label{eq:HB-k}
\mu(n)=\sum_{j=1}^{k}(-1)^{j-1}\binom{k}{j}
\bigl(\mu_u^{*j}*1^{*(j-1)}\bigr)(n),
\end{equation}
because $(\mu-\mu_u)^{*k}*1^{*(k-1)}$ vanishes on $n\le u^k$; expand by the
binomial theorem and use $\mu*1=\delta$. Thus determination by the prefix
is not special to $u^2$. What is special to $t\le u^2$ is the order-two form
\eqref{eq:bilinear-excursion}, with rank-one coefficients $\mu(d)\mu(j)$ and
elementary kernel $\lfloor t/(dj)\rfloor$. This is a truncation identity of
Heath-Brown/Vaughan type \cite{HB82}; summed forms of this kind are used in
the fast computation of $M$ \cite{HT23}.
\end{remark}

\begin{question}[Square-excursion bridges]
\label{q:excursion-constraint}
Let $u=x_{0,\nu}$ and $x_{1,\nu}\le u^2$. Then
$\mathcal B_u(u)=\mathcal B_u(x_{1,\nu})=0$ and
$-\sigma_\nu\mathcal B_u(t)>0$ for $u<t<x_{1,\nu}$. Can the rank-one
coefficients $\mu(d)\mu(j)$, the kernel $\lfloor t/(dj)\rfloor$ and these
bridge and sign constraints be combined to prove
\[
\max_{u<t\le x_{1,\nu}}|\mathcal B_u(t)|
\ll_\eps x_{1,\nu}^{1/2+\eps}
\]
uniformly over square-ordinary excursions? By
Corollary~\ref{cor:power-ordinary-RH} this uniform bound is equivalent to
RH. By Corollary~\ref{cor:linear-window-RH} it suffices to consider
$t\le\min\{x_{1,\nu},\frac54u\}$, where $\mathcal B_u(t)$ reduces to
the localized linear form and uses only $\mu(d)$ with $d\le5u/8$.
\end{question}

\subsection{Mellin mass of a complete excursion}

\begin{proposition}[Excursion blocks of $1/\zeta$]
\label{prop:excursion-mellin}
For a complete excursion $I=(x_0,x_1]$ put
$D_I(s)=\sum_{x_0<n\le x_1}\mu(n)n^{-s}$. Then
\begin{equation}
\label{eq:excursion-abel}
D_I(s)=\sum_{n=x_0+1}^{x_1-1}M(n)
\bigl(n^{-s}-(n+1)^{-s}\bigr).
\end{equation}
For real $\varsigma>0$, $D_I(\varsigma)$ has the sign of the excursion and
\begin{equation}
\label{eq:excursion-mass-lower}
|D_I(\varsigma)|\ge
\varsigma\beta_I x_1^{-\varsigma-1}
\ge\varsigma H_I^2x_1^{-\varsigma-1}.
\end{equation}
For $\Re s>1$,
\begin{equation}
\label{eq:1overzeta-excursions}
\frac1{\zeta(s)}=\sum_I D_I(s),
\end{equation}
zero runs contributing nothing.
\end{proposition}

\begin{proof}
Summation by parts gives \eqref{eq:excursion-abel} because $M$ vanishes at
both endpoints. All interior $M(n)$ have the sign of the excursion, and
$n^{-\varsigma}-(n+1)^{-\varsigma}\ge\varsigma x_1^{-\varsigma-1}$ for $n<x_1$. The
bound $\beta_I\ge H_I^2$ is \eqref{eq:beta-LH}. If $M$ vanishes at two
consecutive integers the intervening $\mu(n)$ is zero, so every $n$ with
$\mu(n)\ne0$ lies in an excursion; sum the disjoint blocks in the
half-plane of absolute convergence.
\end{proof}

\section{Trial-division horizon of an excursion}
\label{sec:prime_horizon}

\begin{lemma}[Square horizon]
\label{lem:square-horizon}
Let $u\ge2$. Trial division by the primes $p\le u$ alone determines $\mu(n)$
for every $n\le u^2$.
\end{lemma}

\begin{proof}
Remove the prime factors $p\le u$ of $n$; if one occurs twice, $\mu(n)=0$.
The cofactor $q\le u^2$ has all prime factors $>u$, so it is $1$ or a single
prime.
\end{proof}

\begin{remark}[Determinacy is not information]
\label{rem:information}
\label{warn:information_language}
Primes $>u$ are not logically new information, and determination by a
finite prefix is not special to $u^2$: \eqref{eq:HB-k} propagates a prefix
to every power range, and \eqref{eq:hyperbola-M} determines $M$ from $M(1)$.
The ranges $(u,2u]$ and $(u,u^2]$ are special only in the concrete sense
that the exact summatory formula is of order one
(\eqref{eq:linear-window-M}) or two (\eqref{eq:bilinear-excursion}) in the
prefix, and that one-stage trial division by primes $\le u$ suffices on
$(u,u^2]$.
\end{remark}

\begin{definition}[Smooth and rough parts]
\label{def:smooth_rough}
For $y>1$ write $n=s_y(n)\,r_y(n)$, where every prime factor of $s_y(n)$
is $\le y$ and every prime factor of $r_y(n)$ is $>y$. Let $\Pp(n)$ be the
largest prime factor of $n$, with $\Pp(1)=1$.
\end{definition}

\begin{lemma}
\label{lem:rough_part}
$\mu(n)=\mu(s_y(n))\,\mu(r_y(n))$. If $n<y^k$ then $\Omega(r_y(n))\le k-1$,
and if $n\le y^2$ then $r_y(n)$ is $1$ or a prime.
\end{lemma}

\begin{proof}
The two parts are coprime. If $\Omega(r)\ge1$ then $y^{\Omega(r)}<r\le n$.
\end{proof}

\begin{remark}[The horizon $u^2$ is not the multiplicative threshold of an excursion]
\label{rem:horizon_inactive}
For a square-ordinary excursion $\sqrt{v_\nu}\le u_\nu$. Every $n\le v_\nu$
has at most one prime factor $>\sqrt{v_\nu}$, so the threshold that governs
the multiplicative structure inside $I_\nu$ is $\sqrt{v_\nu}$, below
$u_\nu$. The regime $t>u_\nu^2$ of Section~\ref{sec:prime_copy_propagation}
lies beyond $v_\nu$ for a square-ordinary excursion; that section concerns
$M(t)$ globally, and Section~\ref{sec:local_large_prime} gives the local
version inside an excursion.
\end{remark}

\section{Large-prime decomposition at level \texorpdfstring{$U^2$}{U^2}}
\label{sec:prime_copy_propagation}

The identities of this section are decompositions by large prime factors of
Buchstab/Legendre type. The asymptotic statements use only the prime number
theorem with classical error term and the Vinogradov--Korobov bound
\cite{MontgomeryVaughan2007,Walfisz1963}.

\begin{definition}[Smooth baseline and large-prime part]
\label{def:baseline_channel}
For $U\ge2$ and $t\ge1$ put
\begin{equation}
\label{eq:baseline_channel_def}
  M^{\rm sm}_U(t):=\sum_{\substack{n\le t\\ \Pp(n)\le U^2}}\mu(n),\qquad
  M^{\rm lp}_U(t):=\sum_{\substack{n\le t\\ \Pp(n)>U^2}}\mu(n),\qquad
  M(t)=M^{\rm sm}_U(t)+M^{\rm lp}_U(t).
\end{equation}
\end{definition}

\begin{proposition}[Single-prime copy]
\label{prop:single_prime_copy}
For a prime $p$ and $t<p^2$,
\begin{equation}
\label{eq:single-prime-copy}
  \Delta_p(t):=\sum_{a\le t/p}\mu(ap)=-M(t/p).
\end{equation}
\end{proposition}

\begin{proof}
$a\le t/p<p$, so $p\nmid a$ and $\mu(ap)=-\mu(a)$.
\end{proof}

\begin{proposition}[Sign content of one prime cone]
\label{prop:prime-cone-signs}
Let $\Sq(y)=\sum_{a\le y}\mu(a)^2$. Among the numbers $ap$ with $a\le y<p$,
the numbers of values $+1,-1,0$ of $\mu$ are
\begin{equation}
\label{eq:cone_counts}
  N_+=\frac{\Sq(y)-M(y)}{2},\qquad N_-=\frac{\Sq(y)+M(y)}{2},\qquad
  N_0=\lfloor y\rfloor-\Sq(y),
\end{equation}
with limiting proportions $3/\pi^2,\ 3/\pi^2,\ 1-6/\pi^2$.
\end{proposition}

\begin{proof}
$\mu(ap)=-\mu(a)$, and $\#\{a\le y:\mu(a)=\mp1\}=(\Sq(y)\mp M(y))/2$.
\end{proof}

\begin{proposition}[First-order window]
\label{prop:linear-new-primes}
If $U^2<t<U^4$, no $n\le t$ has two prime factors $>U^2$ (with
multiplicity), and
\begin{equation}
\label{eq:first-order-total}
  M^{\rm lp}_U(t)=-\sum_{U^2<p\le t}M(t/p).
\end{equation}
\end{proposition}

\begin{proof}
Two such factors would give $n>U^4>t$. Hence $M^{\rm lp}_U(t)=\sum_{U^2<p\le t}
\Delta_p(t)$, and $t<U^4<p^2$; apply \eqref{eq:single-prime-copy}.
\end{proof}

\begin{proposition}[Prime-count form]
\label{prop:prime_count_form}
For $U^2<t<U^4$ and $R:=t/U^2$,
\begin{equation}
\label{eq:prime_count_form}
  M^{\rm lp}_U(t)=-\sum_{a\le R}\mu(a)\bigl(\pi(t/a)-\pi(U^2)\bigr),
  \qquad
  M(t)=M^{\rm sm}_U(t)-\sum_{a\le R}\mu(a)\pi(t/a)+M(R)\pi(U^2).
\end{equation}
\end{proposition}

\begin{proof}
Expand \eqref{eq:first-order-total} as $-\sum_{p}\sum_{a\le t/p}\mu(a)$ and
interchange: for fixed $a$ the prime runs over $U^2<p\le t/a$, which is
non-empty only if $a<R$.
\end{proof}

Thus the only input from primes $>U^2$ is the value of $\pi$ at the points
$t/a$, $a\le R$, weighted by $\mu(a)$ at the small scale $R<U^2$.

\begin{definition}[Copy multiplicities]
\label{def:copy-multiplicity}
For $U^2<t<U^4$, $R=t/U^2$ and $1\le y\le\lfloor R\rfloor$ put
\begin{equation}
\label{eq:Ny-exact}
  N_y(t,U)=\#\Bigl\{p:\ \max\bigl(U^2,\tfrac{t}{y+1}\bigr)<p\le\tfrac ty\Bigr\},
  \qquad\text{so that}\qquad
  M^{\rm lp}_U(t)=-\sum_{1\le y\le\lfloor R\rfloor}N_y(t,U)M(y).
\end{equation}
\end{definition}

\begin{remark}[Asymptotics of the multiplicities]
\label{rem:copy-count-asymptotic}
For fixed $R>1$ and $U\to\infty$,
\begin{equation}
\label{eq:Ny-fixed-R}
  N_y\sim\frac{Ru^2}{2\log U}\cdot\frac1{y(y+1)}\quad(1\le y\le\lfloor R\rfloor-1),
  \qquad
  N_{\lfloor R\rfloor}\sim\frac{U^2}{2\log U}\Bigl(\frac R{\lfloor R\rfloor}-1\Bigr)
  \quad(R\notin\mathbb{Z}),
\end{equation}
and $N_R=0$ if $R\in\mathbb{Z}$. The top multiplicity is truncated because
$t/(\lfloor R\rfloor+1)<U^2$.
\end{remark}

\begin{definition}[Coherent coefficient]
\label{def:coherent_coefficient}
\begin{equation}
\label{eq:C_def}
  \Cc(R):=\int_1^R\frac{M(w)}{w^2}\,dw
  =\sum_{a\le R}\mu(a)\Bigl(\frac1a-\frac1R\Bigr)
  =\sum_{a\le R}\frac{\mu(a)}a-\frac{M(R)}R .
\end{equation}
\end{definition}

\begin{proposition}[Fixed $R$]
\label{prop:fixed_R_corrected}
For fixed $R>1$, as $U\to\infty$,
\begin{equation}
\label{eq:E_fixed_R}
  M^{\rm lp}_U(Ru^2)=-\frac{Ru^2}{2\log U}\bigl(\Cc(R)+o(1)\bigr).
\end{equation}
\end{proposition}

\begin{proof}
In \eqref{eq:prime_count_form} the sum is finite; for $a<R$ the prime number
theorem gives $\pi(Ru^2/a)-\pi(U^2)\sim(R/a-1)U^2/(2\log U)$, and
$\sum_{a\le R}\mu(a)(R/a-1)=R\,\Cc(R)$.
\end{proof}

\begin{lemma}[The multiplicity-weighted coefficient]
\label{lem:copy-weight-coefficient}
The coefficient
\begin{equation}
\label{eq:Theta_def}
  \Upsilon(R):=\sum_{1\le y\le R}\frac{M(y)}{y(y+1)}
  =\sum_{n\le R}\frac{\mu(n)}n-\frac{M(R)}{\lfloor R\rfloor+1}
\end{equation}
satisfies
\begin{equation}
\label{eq:Theta_C}
  \Upsilon(R)=\Cc\bigl(\lfloor R\rfloor+1\bigr),\qquad
  \Cc(R)=\Upsilon(R)-M(R)\Bigl(\frac1R-\frac1{\lfloor R\rfloor+1}\Bigr),
\end{equation}
and
\begin{equation}
\label{eq:C_tail}
  \Cc(R)=-\int_R^\infty\frac{M(w)}{w^2}\,dw,\qquad
  \Upsilon(R)=-\sum_{n>R}\frac{M(n)}{n(n+1)}.
\end{equation}
\end{lemma}

\begin{proof}
$M$ is constant on $[y,y+1)$ and $\int_y^{y+1}w^{-2}dw=1/(y(y+1))$, which
gives \eqref{eq:Theta_C}. For $s>1$, \eqref{eq:mellin_M} gives
$\int_1^\infty M(w)w^{-s-1}\,dw=1/(s\zeta(s))\to0$ as $s\downarrow1$. For
$s\in[1,1+\delta]$ the bound $M(w)\ll we^{-c\sqrt{\log w}}$ dominates the
integrand by an integrable multiple of $e^{-c\sqrt{\log w}}/w$, so dominated
convergence gives $\int_1^\infty M(w)w^{-2}dw=0$.
\end{proof}

\begin{remark}[$\Cc$ versus $\Theta$]
\label{rem:C-vs-Theta}
Summing \eqref{eq:Ny-fixed-R} against $M(y)$ suggests the coefficient
$\Upsilon(R)$, but the truncated top multiplicity replaces it by $\Cc(R)$; by
\eqref{eq:Theta_C} the two agree exactly when $M(\lfloor R\rfloor)=0$. For
$1<R<2$ one has $M^{\rm lp}_U=-(\pi(t)-\pi(U^2))\sim-(R-1)U^2/(2\log U)$, i.e.\
$\Cc(R)=1-1/R$, whereas $\Upsilon(R)=1/2$.
\end{remark}

\begin{corollary}[Decay of the coefficients]
\label{cor:AR-VK}
By \eqref{eq:C_tail} and \eqref{eq:VK_M},
\begin{equation}
\label{eq:AR-VK}
  \Cc(R),\ \Upsilon(R)\ll\exp\bigl(-c_1(\log R)^{3/5}(\log\log R)^{-1/5}\bigr).
\end{equation}
\end{corollary}

\begin{proposition}[Both parts are large]
\label{prop:both_parts_large}
\begin{enumerate}[label=\textup{(\roman*)}]
\item For fixed $R>1$ with $\Cc(R)\ne0$,
\begin{equation}
\label{eq:E_fixed_R_size}
  M^{\rm lp}_U(Ru^2)\sim-\frac{Ru^2}{2\log U}\Cc(R),\qquad
  M^{\rm sm}_U(Ru^2)\sim+\frac{Ru^2}{2\log U}\Cc(R);
\end{equation}
for $1<R<2$, $\Cc(R)=1-1/R>0$.
\item There is $C>0$ such that, as $t\to\infty$, uniformly for $U^2<t<U^4$
with $R=t/U^2\ge\exp(C(\log\log t)^2)$,
\begin{equation}
\label{eq:E_log2}
  M^{\rm lp}_U(t)=(1+o(1))\frac{t}{(\log t)^2},\qquad
  M^{\rm sm}_U(t)=-(1+o(1))\frac{t}{(\log t)^2}.
\end{equation}
\end{enumerate}
In particular, in both regimes $|M^{\rm lp}_U(t)|$ and $|M^{\rm sm}_U(t)|$ exceed
$t^{1/2+\eps}$ for every fixed $\eps<1/2$ and all large $t$, and only their
sum $M$ can be compared with the scale $t^{1/2+\eps}$.
\end{proposition}

\begin{proof}
We use the classical form \cite{MontgomeryVaughan2007}: for some $c>0$,
\begin{equation}
\label{eq:pnt_inputs}
  \pi(y)=\li(y)+O\bigl(ye^{-c\sqrt{\log y}}\bigr),\quad
  M(y)\ll ye^{-c\sqrt{\log y}},\quad
  \mt(y):=\sum_{a\le y}\frac{\mu(a)}a\ll e^{-c\sqrt{\log y}},
\end{equation}
and $\sum_{n\ge1}\mu(n)\log n/n=-1$ as a convergent series (the value at
$s=1$ of $-\frac{d}{ds}\zeta(s)^{-1}$, since $\zeta'/\zeta^2\to-1$). The
integral below converges absolutely by the bound on $\mt$, and
\begin{equation}
\label{eq:mtilde_integral}
  \int_1^\infty\mt(w)\frac{dw}w
  =\lim_{Y\to\infty}\Bigl(\mt(Y)\log Y-\sum_{a\le Y}\frac{\mu(a)\log a}a\Bigr)=1.
\end{equation}

(i) Proposition~\ref{prop:fixed_R_corrected}, and $M(t)=o(t/\log t)$. For
$1<R<2$, $M=1$ on $[1,2)$.

(ii) Since $t<U^4$, $R<\sqrt t$, so $\log(t/a)\ge\frac12\log t$ for $a\le R$.
By \eqref{eq:prime_count_form},
\begin{equation}
\label{eq:E_three_terms}
  M^{\rm lp}_U(t)=-\sum_{a\le R}\mu(a)\li(t/a)
  -\sum_{a\le R}\mu(a)\bigl(\pi(t/a)-\li(t/a)\bigr)+M(R)\pi(U^2).
\end{equation}
The second term is $\ll t(1+\log R)e^{-c\sqrt{\log t/2}}=o(t/\log^2t)$. For
the third, $\tfrac14\log t<\log U<\tfrac12\log t$, so
\[
|M(R)|\pi(U^2)\ll Re^{-c\sqrt{\log R}}\frac{U^2}{\log U}
\ll e^{-c\sqrt{\log R}}\frac{t}{\log t}=o(t/\log^2t)
\]
once $c\sqrt C>1$. For the first, with $F(w)=w\li(t/w)$,
\begin{align}
\label{eq:partial_summation_li}
  \sum_{a\le R}\mu(a)\li(t/a)
  &=\mt(R)F(R)-\int_1^R\mt(w)F'(w)\,dw,\\
  F'(w)&=\li(t/w)-\frac{t/w}{\log(t/w)}
  =\frac{t/w}{\log^2(t/w)}\Bigl(1+O\Bigl(\frac1{\log t}\Bigr)\Bigr).\notag
\end{align}
The boundary term is $\ll e^{-c\sqrt{\log R}}t/\log t=o(t/\log^2t)$. The
integral equals
\[
\frac t{\log^2t}\left(\int_1^R\mt(w)g_t(w)\frac{dw}w+O(1/\log t)\right),
\qquad
g_t(w)=\frac{\log^2t}{\log^2(t/w)}\in[1,4].
\]
Extend $g_t(w)\mathbf1_{\{w\le R\}}$ by zero to $[1,\infty)$. Since
$R\ge\exp(C(\log\log t)^2)\to\infty$, this factor tends to $1$ pointwise and is
bounded by $4$, while $|\mt(w)|/w$ is integrable; dominated convergence gives
$\int_1^R\mt(w)g_t(w)\,dw/w=1+o(1)$, where the $o(1)$ depends only on $t$
and tends to $0$ as $t\to\infty$ uniformly over all pairs $(U,R)$ in the
stated range: indeed $g_t(w)\mathbf 1_{\{w\le R\}}\to1$ uniformly for $w$ in
compact sets, since $R\ge\exp(C(\log\log t)^2)$ and $w\le W$ give
$|g_t(w)-1|\ll\log W/\log t$. Hence
$M^{\rm lp}_U(t)=(1+o(1))t/\log^2t$, and $M^{\rm sm}_U=M-M^{\rm lp}_U$ with $M(t)=o(t/\log^2t)$.
\end{proof}

\begin{remark}[Why $t/\log^2t$]
\label{rem:why_log2}
$\Cc(R)$ captures only the $1/\log t$ term, whose coefficient tends to zero
by Corollary~\ref{cor:AR-VK}. The $1/\log^2t$ term comes from
$\zeta'(s)/\zeta(s)^2\to-1$, that is, from the pole of $\zeta$ at $s=1$,
and does not decay in $R$.
\end{remark}

\begin{remark}[First copy transitions]
\label{rem:first-copy-transitions}
For $1<R<2$ every copy has length $1$ and $M^{\rm lp}_U(t)=-(\pi(t)-\pi(U^2))$. For
$2<R<3$ copies of length $2$ appear but $M(2)=0$. For $3<R<4$ copies of
length $3$ enter with $M(3)=-1$, of opposite net sign.
\end{remark}

\section{Local large-prime decomposition of an excursion}
\label{sec:local_large_prime}

\begin{proposition}[Local decomposition]
\label{prop:local_large_prime}
Let $I_\nu=(u,v]$ be a complete excursion and $y=\sqrt v$. For $u<t\le v$,
\begin{equation}
\label{eq:local_split}
  M(t)=\sum_{\substack{u<n\le t\\ \Pp(n)\le y}}\mu(n)
  -\sum_{a<y}\mu(a)\,\Pi_a(t),\qquad
  \Pi_a(t):=\#\Bigl\{p:\ \max\bigl(y,\tfrac ua\bigr)<p\le\tfrac ta\Bigr\},
\end{equation}
and the bridge condition $M(v)=0$ reads
\begin{equation}
\label{eq:local_bridge}
  \sum_{\substack{u<n\le v\\ \Pp(n)\le\sqrt v}}\mu(n)
  =\sum_{a<\sqrt v}\mu(a)\,\Pi_a(v).
\end{equation}
\end{proposition}

\begin{proof}
$M(u)=0$. An $n\le v=y^2$ has at most one prime factor $p>y$; then
$n=ap$ with $a<v/y=y<p$, so $\mu(n)=-\mu(a)$, and $(a,p)\mapsto ap$ is a
bijection onto these $n$.
\end{proof}

\begin{remark}[Primes in short intervals]
\label{rem:local_meaning}
The decomposition holds for every excursion, with threshold $\sqrt{v_\nu}$.
The large-prime input consists of prime counts in the intervals
$(u/a,t/a]$, of length at most $\ell_\nu/a$ at height about $v_\nu/a$.
Huxley proved the prime number theorem uniformly in intervals of length
$h\ge x^{7/12+\eps}$ \cite{Huxley1972}; the zero-density estimates of
Guth--Maynard give intervals of length $x^{17/30+o(1)}$
\cite{GuthMaynard2024}. For M\"obius sums, Matom\"aki--Ter\"av\"ainen prove
$\sum_{x<n\le x+x^\theta}\mu(n)=o(x^\theta)$ for every $\theta>0.55$
\cite{MatomakiTeravainen2023}. The cited results do not provide, below these ranges, an asymptotic for
$\Pi_a$ uniform in $a$.
\end{remark}

\begin{lemma}[Inverse logarithms]
\label{lem:inverse_log_sum}
Let $C$ be as in Proposition~\ref{prop:both_parts_large}. Uniformly for
\[
\exp\bigl(C(\log\log x)^2\bigr)\le z\le\sqrt x,
\]
\begin{equation}
\label{eq:inverse_log_sum}
  \sum_{a\le z}\frac{\mu(a)}{a\log(x/a)}=-\frac{1+o(1)}{(\log x)^2}.
\end{equation}
\end{lemma}

\begin{proof}
Partial summation gives
\[
\frac{\mt(z)}{\log(x/z)}-
\int_1^z\frac{\mt(w)}{w\log^2(x/w)}\,dw.
\]
Since $z\le\sqrt x$, the boundary term is
$O(e^{-c\sqrt{\log z}}/\log x)=o(1/\log^2x)$ uniformly. Write
$\log^{-2}(x/w)=g_x(w)\log^{-2}x$ with $1\le g_x(w)\le4$, extend
$g_x\mathbf1_{\{w\le z\}}$ by zero, and apply dominated convergence with
\eqref{eq:mtilde_integral}.
\end{proof}

\begin{heuristic}[Opposite drifts]
\label{heu:local_drift}
Let $u\ge\sqrt v\exp(C(\log\log v)^2)$. Near $t\in I_\nu$ the counts
$\Pi_a$ grow for $a<t/\sqrt v$, at rate $1/(a\log(t/a))$ if primes near
$t/a$ have their expected local density uniformly in $a$. Over $(t,t+h]$
the large-prime part then changes by about
$-h\sum_{a<t/\sqrt v}\mu(a)/(a\log(t/a))=(1+o(1))h/\log^2t$ by
Lemma~\ref{lem:inverse_log_sum}. Since $M(v)-M(u)=0$, the $\sqrt v$-smooth
part drifts at rate about $-1/\log^2v$, independently of $\sigma_\nu$. As
\[
\frac{\ell_\nu}{\log^2v_\nu}\gg\sqrt{\ell_\nu}
\iff
\ell_\nu\gg\log^4v_\nu,
\]
for $\ell_\nu\gg(\log v_\nu)^4$ both drift contributions exceed
$\sqrt{\ell_\nu}$, and \eqref{eq:local_bridge} is a balance between two
opposite drifts: the local analogue of
Proposition~\ref{prop:both_parts_large}.
\end{heuristic}

\section{A one-frequency model for the black curve}
\label{sec:model}

\begin{heuristic}[Scale and period]
\label{heu:model}
Motivated by \eqref{eq:explicitM}, model $M(x)\approx\sqrt x\,F(\log x)$
with $F$ almost periodic, and retain only the lowest zero $\gamma_1$ (RH
does not imply that its coefficient dominates). Within an excursion of
length $\asymp x$ this gives $H_\nu\asymp\sqrt x$ and
$\sqrt{H_\nu}\ell_\nu\asymp x^{5/4}$, so the black curve oscillates at scale
\begin{equation}
\label{eq:model_scale}
  |\Anet(x)|\asymp x^{5/4},
\end{equation}
the scale of Corollary~\ref{cor:rh}, with period $2\pi/\gamma_1$ in
$\log x$, i.e.\ a factor $e^{2\pi/\gamma_1}\approx1.5597$ in $x$. In this
model consecutive zeros of $M$ are separated by the factor
$e^{\pi/\gamma_1}=\lambda_*$, the threshold of
Corollary~\ref{cor:descending}(iii): a model excursion is exactly
$\lambda_*$-long.
\end{heuristic}

\section{Numerical data}
\label{sec:data}

The following census covers $n\le N=455\,052\,511$, with $Z(N)=87\,385$
excursions, recording $(u_\nu,\tau_\nu,v_\nu,H_\nu,\sigma_\nu)$ and the
curves $\Ap,\Am,\Anet$:
\begin{equation}
\label{eq:census_values}
  \Ap(N)=6.3449\times10^9,\quad \Am(N)=-7.0877\times10^9,\quad
  \Anet(N)=-7.4279\times10^8,
\end{equation}
\begin{equation}
\label{eq:census_max}
  \max_\nu\frac{H_\nu}{\sqrt{\ell_\nu}}=2.0743,\qquad
  \max_\nu\frac{\beta_\nu}{\ell_\nu^{3/2}}=1.0214 .
\end{equation}
Since $|M(x)|\le0.571\sqrt x$ for $33\le x\le10^{16}$ \cite{Hurst18}, and
$|M(n)|\le\sqrt n$ for $n<33$ with equality only at $n=1$, the
constant $C_0=1$ is admissible in Lemma~\ref{lem:transfer}, which gives
$\Ap(N)+|\Am(N)|\le\frac45N^{5/4}=5.317\times10^{10}$ against the observed
$1.3433\times10^{10}$, a ratio of $3.96$. Heuristic~\ref{heu:renewal} with
$L=\sqrt N$ predicts $\#\{\ell_\nu>\sqrt N\}\asymp N^{1/4}\approx146$. At
this range $N^{0.05}\approx2.7$, so factors $x^\eps$ are not separable from
constants. Reproducing these values requires the excursion file, the
area-integration convention of Proposition~\ref{prop:explicit}, and the
terminal convention of Definition~\ref{def:Ms-terminal}. These data are not used in any proof.

\section{Open problems}
\label{sec:open_questions}

\begin{enumerate}[label=\textup{(\arabic*)}]
\item In the localized form of Corollary~\ref{cor:linear-window-RH}, bound
\[
\sum_{d\le5u/8}\mu(d)
\bigl(\lfloor t/d\rfloor-\lfloor u/d\rfloor\bigr)
\]
by $t^{1/2+\eps}$ on
$u<t\le\min\{x_1,5u/4\}$, using the full simultaneous intra-excursion
constraints rather than pointwise estimates.
\item Quantify the multiscale persistence forced by
Lemma~\ref{lem:channel-persistence}.  By Corollary~\ref{cor:channel-descent}
and Remark~\ref{rem:channel-loss}, channel descent iterates but loses the
factor $2K/k^\alpha$ per step; the problem is to show that along a descent
chain from an RH-violating value the loss is eventually absent, or that the
deep remainder $\mathcal R_K$ carries the value and can be bounded.
\item Question~\ref{q:layer-depth} in the form left by
Corollary~\ref{cor:bounded-depth}: bounded depth fixes only $o(u)$ window
increments, whose sum is of order $u(\log\log u)^k/\log u$; decide whether
the unit-step and excursion constraints on the remaining free coordinates
$\delta_d$, $\Omega(d)>k$, already force the cancellation down to
$u^{1/2+\eps}$. The first case is $k=1$ (primes fixed, $\delta_p=-1$).
\item Can the threshold $\lambda_*=e^{\pi/\gamma_1}$ of
Corollary~\ref{cor:descending} be lowered? By
Remark~\ref{rem:lambda-star} this requires information on the position of
large values inside excursions, beyond the count
\eqref{eq:KP-large-excursions}.
\item Whether \eqref{eq:rh_area} is equivalent to RH.
\item Whether the pruned local condition \eqref{eq:pruned_Q} holds, and
whether the bridge constraints \eqref{eq:Q_def}, \eqref{eq:zeroth-corr}
admit a third independent relation involving the defect
$\ell_\nu/2-H_\nu\ge0$.
\item Establish (SG) for $Z=H/\sqrt\ell$ in the simple random walk model of
the M\"obius walk, and compute the extreme scale of
Corollary~\ref{cor:ev-square} there; the Brownian-excursion limit suggests
$c=2$.
\item Numerically: the number of excursions with $\ell_\nu>\sqrt N$ against
$N^{1/4}$, the maxima of $H_\nu/\sqrt{\ell_\nu}$ and $\beta_\nu/\ell_\nu^{3/2}$
over them; the split \eqref{eq:local_split} and the $\sigma_\nu$-independence
of Heuristic~\ref{heu:local_drift}; the frequency of touching runs and of
ballistic excursions (Heuristic~\ref{heu:local_fails}); and the distribution
of $x_1/x_0$, to be compared with $\lambda_*$.
\end{enumerate}

\begin{thebibliography}{99}

\bibitem{CP17}
P.~Carmona and N.~P\'etr\'elis, \emph{Limit theorems for random walk excursion
conditioned to enclose a typical area}, arXiv:1709.06448.

\bibitem{DKW15}
D.~Denisov, M.~Kolb and V.~Wachtel, \emph{Local asymptotics for the area of
random walk excursions}, J. London Math. Soc. \textbf{91} (2015), 495--513.

\bibitem{GuthMaynard2024}
L.~Guth and J.~Maynard, \emph{New large value estimates for Dirichlet
polynomials}, Ann. of Math. (2) \textbf{203} (2026), no.~2, 623--675,
doi:10.4007/annals.2026.203.2.6.

\bibitem{HB82}
D.~R. Heath-Brown, \emph{Prime numbers in short intervals and a generalized
Vaughan identity}, Canad. J. Math. \textbf{34} (1982), 1365--1377.

\bibitem{HT23}
H.~A. Helfgott and L.~Thompson, \emph{Summing $\mu(n)$: a faster elementary
algorithm}, Res. Number Theory \textbf{9} (2023), Paper No.~6.

\bibitem{Hurst18}
G.~Hurst, \emph{Computations of the Mertens function and improved bounds on
the Mertens conjecture}, Math. Comp. \textbf{87} (2018), no.~310, 1013--1028.

\bibitem{Huxley1972}
M.~N. Huxley, \emph{On the difference between consecutive primes}, Invent.
Math. \textbf{15} (1972), 164--170.

\bibitem{Janson07}
S.~Janson, \emph{Brownian excursion area, Wright's constants in graph
enumeration, and other Brownian areas}, Probab. Surveys \textbf{4} (2007),
80--145.

\bibitem{KP86}
J.~Kaczorowski and J.~Pintz, \emph{Oscillatory properties of arithmetical
functions. I}, Acta Math. Hungar. \textbf{48} (1986), 173--185.

\bibitem{KotnikTeRiele2006}
T.~Kotnik and H.~te Riele, \emph{The Mertens conjecture revisited}, in
Algorithmic Number Theory (ANTS VII), Lecture Notes in Comput. Sci.
\textbf{4076}, Springer, Berlin, 2006.

\bibitem{MatomakiTeravainen2023}
K.~Matom\"aki and J.~Ter\"av\"ainen, \emph{On the M\"obius function in all
short intervals}, J. Eur. Math. Soc. \textbf{25} (2023), no.~4, 1207--1225.

\bibitem{MontgomeryVaughan2007}
H.~L. Montgomery and R.~C. Vaughan, \emph{Multiplicative Number Theory I:
Classical Theory}, Cambridge Stud. Adv. Math. \textbf{97}, Cambridge
University Press, 2007.

\bibitem{Ng2004}
N.~Ng, \emph{The distribution of the summatory function of the M\"obius
function}, Proc. London Math. Soc. (3) \textbf{89} (2004), no.~2, 361--389.

\bibitem{NguyenThe04}
M.~Nguy\^en Th\^e, \emph{Area and inertial moment of Dyck paths}, Combin.
Probab. Comput. \textbf{13} (2004), 697--716.

\bibitem{OdlyzkoTeRiele1985}
A.~M. Odlyzko and H.~J.~J. te Riele, \emph{Disproof of the Mertens
conjecture}, J. Reine Angew. Math. \textbf{357} (1985), 138--160.

\bibitem{Richard05}
C.~Richard, \emph{On $q$-functional equations and excursion moments},
arXiv:math/0503198.

\bibitem{Takacs91}
L.~Tak\'acs, \emph{A Bernoulli excursion and its various applications},
Adv. Appl. Probab. \textbf{23} (1991), 557--585.

\bibitem{Titchmarsh1986}
E.~C. Titchmarsh, \emph{The Theory of the Riemann Zeta-Function}, 2nd ed.,
revised by D.~R. Heath-Brown, Clarendon Press, Oxford, 1986.

\bibitem{Walfisz1963}
A.~Walfisz, \emph{Weylsche Exponentialsummen in der neueren Zahlentheorie},
VEB Deutscher Verlag der Wissenschaften, Berlin, 1963.

\end{thebibliography}

\end{document}
